% =============================================================================
% Analysis of baseline performance clouds.
% Drop into your Overleaf project and % =============================================================================
% Analysis of baseline performance clouds.
% Drop into your Overleaf project and % =============================================================================
% Analysis of baseline performance clouds.
% Drop into your Overleaf project and % =============================================================================
% Analysis of baseline performance clouds.
% Drop into your Overleaf project and \input{cloud_analysis} where needed.
% Requires: amsmath, booktabs, array.
% Figure paths assume Overleaf layout: figures/external baselines/<name>
% =============================================================================

\section{External Baselines}\label{sec:external-baselines}

We benchmark our service-design framework against seven external baselines drawn from three families: (i) partition with continuous approximation (\textbf{Joint-Nom}, \textbf{Joint-DRO}, \textbf{SP-Lit (Carlsson)}); (ii) routing-based (\textbf{Multi-FR (Jacquillat)}, \textbf{TP-Lit (Liu)}); (iii) operational extremes (\textbf{VCC}, \textbf{OD/fully-on-demand}). Each baseline sweeps its own structural hyperparameters internally per regime point and picks the combination minimising the aggregate cost $C_{\mathrm{prov}} + w_r \cdot \bar{u} \cdot \Lambda$. The comparison is stress-tested at three tiers: \S\ref{sec:cloud-analysis} synthetic $5{\times}5$ grid; \S\ref{sec:real-cloud} a 25-block-group Pittsburgh slice with real road distances and ACS population demand; \S\ref{sec:multi-city} four U.S.\ cities spanning topology archetypes (Pittsburgh, Chicago, Manhattan, LA). To make the regime response readable, each cloud is presented in two faceted views: one figure with $\Lambda$ fixed per panel and $w_r$ varying within the panel (shows how each baseline responds to the user-time weight at a given demand level), and a companion figure with $w_r$ fixed per panel and $\Lambda$ varying within the panel (shows how each baseline responds to demand at a given user-time weight).

\subsection{Analysis on the synthetic instance}\label{sec:cloud-analysis}

\begin{figure}[t]
\centering
\includegraphics[width=0.92\linewidth]{figures/external baselines/cloud_synthetic}
\caption{Performance cloud on the $5{\times}5$ synthetic instance ($\Lambda \in \{50, 150, 300\}$\,pax/h, $w_r \in \{1, 3, 10\}$, $w_v{=}1$, $c_f{=}4.0$\,\$/veh-h). Marker size $\propto \Lambda$, fill shade $\propto w_r$ (lighter $=$ larger $w_r$), shaded hull is each baseline's spread over the $3{\times}3$ regime grid.}
\label{fig:cloud}
\end{figure}

\begin{figure}[t]
\centering
\includegraphics[width=\linewidth]{figures/external baselines/cloud_synthetic_by_lambda}
\caption{Synthetic cloud, fixed $\Lambda$ per panel, $w_r$ varying within panel (fill shade encodes $w_r$, lighter $=$ smaller $w_r$). Lines connect the $w_r$ sweep for each baseline. Baselines whose design is genuinely $w_r$-insensitive at a given $\Lambda$ (TP-Lit at every $\Lambda$, VCC and OD at low $\Lambda$) have all three $w_r$ rows at identical (provider, user) coordinates; markers are then jittered slightly in $\log x$ so all three remain visible.}
\label{fig:cloud-by-lambda}
\end{figure}

\begin{figure}[t]
\centering
\includegraphics[width=\linewidth]{figures/external baselines/cloud_synthetic_by_wr}
\caption{Synthetic cloud, fixed $w_r$ per panel, $\Lambda$ varying within panel (marker size $\propto \Lambda$). Lines connect the $\Lambda$ sweep for each baseline.}
\label{fig:cloud-by-wr}
\end{figure}

Figure~\ref{fig:cloud} plots each baseline's provider cost rate (log scale) against average user time per rider across the full $3 \times 3$ regime grid. Figures~\ref{fig:cloud-by-lambda} and~\ref{fig:cloud-by-wr} decompose the same data by holding one regime parameter fixed per panel. The by-$\Lambda$ view exposes each method's $w_r$ trajectory: partition methods (Joint-Nom, Joint-DRO, SP-Lit) sweep from cheap-provider/long-wait at $w_r{=}1$ to expensive-provider/short-wait at $w_r{=}10$, while TP-Lit, VCC, and OD slide only along the provider axis at roughly fixed user time, and Multi-FR barely moves. The by-$w_r$ view exposes each method's $\Lambda$ trajectory: Multi-FR shows the strongest Mohring effect (user time drops from $\approx 1.07$\,h at $\Lambda{=}50$ to $\approx 0.56$\,h at $\Lambda{=}300$, $w_r{=}1$), partition methods are weakly anti-Mohring (Joint's $T^*$ lengthens with $\Lambda$), and OD/VCC have nearly flat user time (the fleet scales up with $\Lambda$ to absorb demand, so all $\Lambda$ variation appears horizontally on the provider axis).

\subsubsection{Per-baseline sweep grids}

\begin{center}
\small
\renewcommand{\arraystretch}{1.15}
\begin{tabular}{lll}
\toprule
Baseline & Swept hyperparameters & Selection \\
\midrule
Joint-Nom  & $K \in \{2,\dots,6\}$ (random-search partition) & best $K$ by aggregate cost \\
Joint-DRO  & $K \in \{2,\dots,6\}$, $\varepsilon = 0.5$ fixed & best $K$ by aggregate cost \\
SP-Lit     & $K \in \{2,\dots,6\}$ $\times$ $T_{\mathrm{fixed}} \in \{0.1, 0.2, 0.3, 0.5, 0.8, 1.0\}$ & best $(K, T)$ \\
Multi-FR   & (fleet, $\delta$) $\in$ $\{(3,0), (5,0), (8,0), (12,0)\}$, internal $T$-grid & best combo \\
TP-Lit     & $\mathtt{max\_fleet} \in \{20, 50, 100\}$; 25 candidate stops & best fleet \\
VCC        & fleet $\in \{5,10,20,30,50\} \times \delta_w \in \{0.1, 0.3, 0.5\}$\,km & best (fleet, $\delta_w$) \\
OD         & fleet $\in \{5,10,15,20,30,50,80\}$ with queueing & best fleet \\
\bottomrule
\end{tabular}
\end{center}

Joint-DRO holds $\varepsilon{=}0.5$ fixed because it is a robustness knob, not a provider/user Pareto knob.

\subsubsection{Joint-Nom and Joint-DRO}

Joint methods occupy the left band of the cloud (provider cost 28 to 136), spanning a wide user-time range (0.11 to 0.78\,h). Two design knobs respond to the regime: the number of districts $K$ (swept externally over $\{2,\dots,6\}$ by the harness, best $K$ picked per regime point by aggregate cost) and the per-district dispatch interval $T^*$ (computed internally by Newton's method for each candidate $K$).

\begin{center}
\small
\begin{tabular}{cccccc}
\toprule
$\Lambda$ & $w_r$ & $K$ & $T^*$ (h) & User time (h) & Provider cost \\
\midrule
 50 &  1 & 4 & 0.88 & 0.60 &  28 \\
 50 & 10 & 5 & 0.13 & 0.12 &  50 \\
150 &  1 & 6 & 0.99 & 0.70 &  38 \\
150 & 10 & 6 & 0.12 & 0.16 &  95 \\
300 &  1 & 6 & 1.00 & 0.78 &  43 \\
300 & 10 & 6 & 0.13 & 0.19 & 124 \\
\bottomrule
\end{tabular}
\end{center}

$T^*$ is highly $w_r$-responsive: at $w_r{=}1$ the optimiser values rider time cheaply and picks $T^* \approx 0.9$ to $1.0$\,h (heavy batching), while at $w_r{=}10$ it drops to $\approx 0.13$\,h (near real-time dispatch) at 3 to 4$\times$ the provider cost. $K$ increases with both $w_r$ and $\Lambda$ (from $K{=}4$ at the cheap end to $K{=}6$ at the expensive end). User time barely responds to $\Lambda$ at low $w_r$ because the Newton $T^*$ formula uses a per-rider wait term $w_r \cdot T$ rather than aggregate $w_r \cdot T \cdot \Lambda$, and the $K$ sweep partially compensates by adding districts at higher $\Lambda$.

Joint-DRO separates from Joint-Nom only at high $\Lambda$ and $w_r$: at $(\Lambda{=}300, w_r{=}10)$, DRO picks fleet 17 / provider cost 136 versus Nom's fleet 15 / provider cost 124 ($\sim 10\%$ premium). Under uniform demand the DRO premium stays modest; non-uniform demand amplifies it.

\subsubsection{Multi-FR (Jacquillat)}

Multi-FR occupies the mid-left region and is the only method with $\Lambda$-decreasing user time on the synthetic grid (the Mohring effect): more riders justify shorter dispatch headways on the same fixed-route backbone, so $T^*$ scales roughly as $1/\sqrt{\Lambda}$. Fleet 3 wins at every synthetic regime point; fleet 5 is available but rarely preferred.

\begin{center}
\small
\begin{tabular}{cccccc}
\toprule
$\Lambda$ & $w_r$ & $T$ (h) & Wait (h) & User time (h) & Provider cost \\
\midrule
 50 &  1 & 1.75 & 0.83 & 1.04 & 30 \\
150 &  1 & 0.81 & 0.40 & 0.62 & 40 \\
300 &  1 & 0.69 & 0.35 & 0.56 & 44 \\
300 & 10 & 0.38 & 0.19 & 0.39 & 69 \\
\bottomrule
\end{tabular}
\end{center}

Multi-FR is one of the few baselines with non-zero walk time ($\approx 0.07$ to $0.12$\,h on the grid), reflecting the rider-to-checkpoint distance. The evaluator assigns each sampled rider to the nearest realised service point globally rather than to a partition-assigned route.

\subsubsection{SP-Lit (Carlsson)}

SP-Lit sits in the centre-right of the cloud, shifted right of Joint despite the same partition-plus-BHH structure. At every regime point its provider cost is 1.6 to 3.1$\times$ Joint-Nom's. The gap is a partition-quality difference: SP-Lit's deterministic ham-sandwich bisection balances the continuous measures $\mu_1$ (spatial) and $\mu_2$ (radial workload) but does not directly minimise per-district cost; Joint's random search with local improvement finds lower-$\alpha_i$ partitions for the same $K$.

\subsubsection{TP-Lit (Liu)}\label{sec:tplit-cloud}

TP-Lit occupies a narrow horizontal band at user time $\approx 0.20$ to $0.23$\,h with provider cost 233 to 244. On the synthetic instance the design is essentially $w_r$-insensitive: the ADP+VRP method ran at its library defaults ($T{=}0.25$\,h, fleet automatically sized by APT to 22), the three $w_r$ rows at each $\Lambda$ are bit-identical, and the three visible markers per $\Lambda$ panel sit on the same coordinate (small log-$x$ jitter applied for visibility). The real-city sweeps (\S\ref{sec:real-cloud}, \S\ref{sec:multi-city}) do add an external $T_{\mathrm{fixed}} \in \{0.15, 0.5\}$\,h $\times$ $\mathtt{max\_fleet} \in \{30, 100\}$ wrapper sweep and pick the best 4-combo by aggregate cost, but the chosen combo is the same at every $w_r$ at most regime points (e.g.\ all three $w_r$ at $\Lambda{=}150$ Pittsburgh pick $T{=}0.5$, $\mathtt{max\_fleet}{=}30$); a $T{=}0.15$ switch only occurs at $w_r{=}10$ on $\Lambda \in \{50, 300\}$. The reason is structural: TP-Lit's inner objective (ADP delivery-time minimisation) does not contain $w_r$, so only the outer wrapper's aggregate-cost selection responds to $w_r$ at all, and that response is quantised by the 4-combo grid. Maintaining a sizeable fleet for 25 candidate stops makes the fleet-cost term alone ($c_f \times \mathtt{fleet}$) account for $\sim 35\%$ of provider cost on the synthetic.

\subsubsection{VCC and OD}

VCC and OD form the operational-extremes panel (Fig.~\ref{fig:cloud-split}b). Both achieve user time below 0.16\,h on the synthetic grid but at provider cost 1 to 2 orders of magnitude above Joint-Nom. VCC sweeps fleet ($\{5,10,20,30,50\}$) and walk tolerance $\delta_w \in \{0.1,0.3,0.5\}$\,km; the optimiser picks larger fleets at higher $\Lambda$ to keep user time constant ($\approx 0.15$\,h), producing a horizontal trajectory in the cloud. OD serves each request individually (depot $\to$ destination $\to$ depot) and is the cost ceiling on provider expense and the cost floor on user time: every other baseline is interpreted relative to OD's extreme. At low $w_r$, OD tolerates moderate queueing (fleet 5--15); at $w_r{=}10$ it upgrades to 50 vehicles at $\Lambda{=}300$.

\subsubsection{Cross-method synthesis}

No single method dominates across the regime grid. The approximate Pareto frontier runs from Joint-Nom/DRO (cheapest provider, highest user time at low $w_r$) through Multi-FR (Mohring-aided at high $\Lambda$) to VCC and OD (lowest user time, highest provider cost). SP-Lit and TP-Lit are dominated on the frontier: SP-Lit by Joint-Nom (same structure, better partitions) and TP-Lit by VCC (lower user time at comparable provider cost). The clearest architectural distinction is the response to $\Lambda$: Multi-FR's user time \emph{decreases} (Mohring); Joint/SP-Lit barely respond (the $K$ sweep absorbs the demand change); VCC/OD are nearly constant in user time (fleet scales with $\Lambda$); TP-Lit is perfectly flat (its design is $\Lambda$-insensitive in this range).

\subsection{Pittsburgh real-data validation}\label{sec:real-cloud}

\begin{figure}[t]
\centering
\includegraphics[width=\linewidth]{figures/external baselines/cloud_pittsburgh}
\caption{Pittsburgh real-data cloud, 25 census block groups in Allegheny County, road-network shortest paths everywhere (optimisation and evaluation), ACS 2023 population-weighted demand. Same regime grid and same axes as Fig.~\ref{fig:cloud}.}
\label{fig:cloud-real}
\end{figure}

\begin{figure}[t]
\centering
\includegraphics[width=\linewidth]{figures/external baselines/cloud_pittsburgh_by_lambda}
\caption{Pittsburgh cloud, fixed $\Lambda$ per panel, $w_r$ varying.}
\label{fig:cloud-real-by-lambda}
\end{figure}

\begin{figure}[t]
\centering
\includegraphics[width=\linewidth]{figures/external baselines/cloud_pittsburgh_by_wr}
\caption{Pittsburgh cloud, fixed $w_r$ per panel, $\Lambda$ varying.}
\label{fig:cloud-real-by-wr}
\end{figure}

We re-run the same regime grid on 25 census block groups around the centroid of HCT's Pittsburgh service area. Demand is ACS-population-weighted ($p_b \in [0.014, 0.062]$ versus the synthetic $p_b{=}1/25$); distances are OSMnx road-network shortest paths (2{,}931 nodes, 7{,}519 edges) everywhere a baseline supports them (Joint-Nom, Joint-DRO, SP-Lit in optimisation; all baselines in evaluation).

\paragraph{Headline numbers at $\Lambda{=}150$, $w_r{=}1$.}
\begin{center}
\small
\begin{tabular}{lcc}
\toprule
Baseline & Provider cost & User time (h) \\
\midrule
Multi-FR (Jacquillat)  &  29.5 & 1.78 \\
Joint-Nom              &  58.6 & 0.92 \\
Joint-DRO              &  59.5 & 0.91 \\
SP-Lit (Carlsson)      & 134.9 & 0.70 \\
TP-Lit (Liu)           & 216.2 & 0.38 \\
VCC                    & 718.1 & 0.21 \\
OD                     & 1220.2 & 0.23 \\
\bottomrule
\end{tabular}
\end{center}

\paragraph{Findings.} Three patterns from the synthetic analysis carry over and one new one emerges. (i) Provider costs scale up across the board because real road distances are longer than Euclidean (typical road-to-Euclidean factor 1.3 to 1.5 in Pittsburgh's grid); OD's provider cost more than triples, while partition methods absorb most of the increase via their road-aware in-district TSPs. (ii) The Pareto ranking is preserved: same ordering Joint $\to$ SP-Lit $\to$ TP-Lit $\to$ VCC $\to$ OD on provider cost, with the reverse on user time. (iii) Joint's $K$ choice tracks regime monotonically ($K{=}4$ at $\Lambda{=}50, w_r{=}1$, $K{=}6$ at $\Lambda{=}300, w_r{=}10$). (iv) Non-uniform demand sharpens the DRO-vs-Nom premium: at $(\Lambda{=}300, w_r{=}10)$, DRO costs $\sim 25\%$ more than Nom (versus $\sim 10\%$ on uniform synthetic), matching the prediction that DRO's worst-case partition diverges more aggressively when demand is spatially non-uniform.

\begin{figure}[t]
\centering
\includegraphics[width=\linewidth]{figures/external baselines/designs_pittsburgh}
\caption{Pittsburgh service designs at $\Lambda{=}150$, $w_r{=}1$, with $K{=}3$ fixed for visualisation (the cloud sweep picks $K$ per regime). Route polylines follow road shortest paths. Partition methods (cols 1 to 3) show contiguous district fills with exact-TSP depot tours. Multi-FR (col 4) shows selected reference loops. TP-Lit (col 5) shows three independent dispatch waves (blues, reds, greens) with realised VRP routes per vehicle. VCC (col 6) shows walk-radius halos ($\delta_w{=}0.3$\,km). OD (col 7) shows depot plus sample on-demand radial trips.}
\label{fig:designs-map}
\end{figure}

The map overlay confirms that the architectural distinctions are also visible in geometry: Joint-Nom and Joint-DRO produce contiguous random-search partitions with a Joint-Nom-optimised depot; SP-Lit's ham-sandwich partition is more axis-aligned; Multi-FR's selected reference loops radiate from the centroid depot; VCC overlays per-block walk halos; OD shows pure depot-anchored trips.

\subsection{Cross-topology validation across four U.S.\ cities}\label{sec:multi-city}

\begin{figure}[t]
\centering
\includegraphics[width=\linewidth]{figures/external baselines/cloud_four_cities}
\caption{Cross-city baseline performance clouds on the shared regime grid. Each panel is one city; markers and shading as in Fig.~\ref{fig:cloud}.}
\label{fig:four-city-cloud}
\end{figure}

\begin{figure}[t]
\centering
\includegraphics[width=\linewidth]{figures/external baselines/cloud_four_cities_by_lambda}
\caption{Cross-city clouds, fixed $\Lambda$ per column, $w_r$ varying within panel. Rows are cities.}
\label{fig:four-city-cloud-by-lambda}
\end{figure}

\begin{figure}[t]
\centering
\includegraphics[width=\linewidth]{figures/external baselines/cloud_four_cities_by_wr}
\caption{Cross-city clouds, fixed $w_r$ per column, $\Lambda$ varying within panel. Rows are cities.}
\label{fig:four-city-cloud-by-wr}
\end{figure}

We re-run the regime sweep on four cities chosen for road-topology contrast: Pittsburgh (irregular hilly, river-cut), Chicago Loop (strict cardinal grid), Manhattan Midtown (avenue-and-street grid), and Los Angeles DTLA (freeway-fragmented grid). Each city uses real TIGER 2023 block groups, ACS-population demand, and full road-network shortest paths in both optimisation and evaluation.

\paragraph{Service-region geometry.} BG count is held at 25 (50 for Manhattan, see below); TIGER block groups are population-sized, so denser cities pack the same count into smaller polygons.

\begin{center}
\small
\begin{tabular}{lcrrrr}
\toprule
City & BGs & Area (km$^2$) & BBox W$\times$H (km) & Diag (km) & Density (pop/km$^2$) \\
\midrule
Pittsburgh & 25 & 24.5 & $7.4 \times 7.2$ & 10.3 &  1{,}003 \\
Chicago    & 25 &  4.6 & $2.3 \times 3.5$ &  4.1 & 13{,}937 \\
Manhattan  & 50 &  2.4 & $2.0 \times 2.5$ &  3.2 & 18{,}891 \\
LA         & 25 &  4.3 & $2.3 \times 2.9$ &  3.7 &  8{,}755 \\
\bottomrule
\end{tabular}
\end{center}

Pittsburgh's 25 BGs cover roughly $10{\times}$ more area than Manhattan's 50, so cross-city cost comparisons reflect both topology and area. Manhattan is enlarged from the default 25 BGs ($1.1$\,km$^2$, too small for FR-Detour's candidate-line library to produce feasible trips) to 50 BGs ($2.4$\,km$^2$); the other three cities are kept at 25 BGs.

\paragraph{Headline numbers at $\Lambda{=}150$, $w_r{=}1$.}
\begin{center}
\small
\begin{tabular}{lrrrr}
\toprule
Baseline & Pittsburgh & Chicago & Manhattan & LA \\
\midrule
\multicolumn{5}{l}{\emph{Provider cost}} \\
Joint-Nom              &   58.6 &   47.6 &   47.9 &   35.5 \\
Joint-DRO              &   59.5 &   45.8 &   54.6 &   36.7 \\
SP-Lit (Carlsson)      &  134.9 &   81.1 &   73.7 &   58.5 \\
Multi-FR (Jacquillat)  &   29.5 &   21.4 &   20.5 &   20.5 \\
TP-Lit (Liu)           &  216.2 &  103.1 &  104.7 &   91.5 \\
VCC                    &  718.1 &  241.9 &  199.7 &  193.3 \\
OD                     & 1220.2 &  476.4 &  417.5 &  343.8 \\
\midrule
\multicolumn{5}{l}{\emph{User time (h per rider)}} \\
Joint-Nom              & 0.92 & 0.64 & 0.58 & 0.56 \\
Joint-DRO              & 0.91 & 0.68 & 0.60 & 0.56 \\
SP-Lit (Carlsson)      & 0.70 & 0.41 & 0.39 & 0.41 \\
Multi-FR (Jacquillat)  & 1.78 & 1.61 & 1.58 & 1.59 \\
TP-Lit (Liu)           & 0.38 & 0.31 & 0.31 & 0.30 \\
VCC                    & 0.21 & 0.12 & 0.11 & 0.09 \\
OD                     & 0.23 & 0.07 & 0.18 & 0.08 \\
\bottomrule
\end{tabular}
\end{center}

\paragraph{Multi-FR calibration.} Multi-FR uses the paper-faithful coverage reward $M = 10{,}000 \cdot w_r$ (Jacquillat et al.\ Table EC.4). With this calibration the master picks fleet $\in \{4, 6, 9, 12\}$ across the regime grid (versus the degenerate fleet $=1$ under our earlier $M = 10 \cdot w_r$), and walk time drops from 0.4 to 0.5\,h down to $\sim 0.07$\,h. However, the master still selects the longest scheduled headway in our $T$-grid ($T_{\max}{=}3$\,h, realised wait $\approx 1.5$\,h) at every regime point because $M{=}10^4 \cdot w_r$ dwarfs the per-passenger wait term $w_r \cdot T/2 \cdot \mathrm{demand}$ by $\sim 6{,}700{\times}$. Multi-FR therefore lands at a many-trips, full-coverage, long-headway design point on real cities. At $w_r{=}10$ on Manhattan and LA the master sometimes flips to an empty-solution corner of the feasible set; see the discussion below.

\paragraph{Cross-method ranking and Pareto positions.} The provider-cost ordering Multi-FR $<$ Joint $<$ SP-Lit $<$ TP-Lit $<$ VCC $<$ OD holds at every $(\Lambda, w_r)$ point on all four cities; the user-time ordering is essentially the reverse. Multi-FR consistently occupies the cheapest-provider corner (provider 20--30, user time 1.6--1.8\,h); OD the opposite (provider 344--1{,}220, user time 0.07--0.23\,h). The Pareto frontier is dense: at $(\Lambda{=}150, w_r{=}1)$, five to seven of the seven baselines are non-dominated on each city. The architectural distinctions between baselines correspond to genuinely different operating points on the same Pareto plane.

\paragraph{Topology controls absolute cost.} Pittsburgh (25 BGs, 24.5\,km$^2$, irregular) inflates absolute provider cost relative to Manhattan (50 BGs, 2.4\,km$^2$, dense grid): Joint-Nom $59$ vs $48$ ($1.2{\times}$), SP-Lit $135$ vs $74$ ($1.8{\times}$), TP-Lit $216$ vs $105$ ($2.1{\times}$), OD $1{,}220$ vs $418$ ($2.9{\times}$). Methods that pay distance per request (OD, VCC) amplify the geometric handicap; partition methods absorb most of it via a shared in-district TSP. Multi-FR is essentially topology-invariant in provider cost (clustering at 20--30 across all four cities) because its long-headway design is dispatch-frequency-bound, not distance-bound. The OD-to-Joint multiplier rises with topology complexity: $9{\times}$ on Manhattan, 10--13$\times$ on Chicago and LA, $21{\times}$ on Pittsburgh.

\paragraph{Regime responsiveness ($w_r$).} As $w_r$ rises from 1 to 10 at $\Lambda{=}150$, partition methods (Joint-Nom/DRO, SP-Lit) drop user time by 66 to 77\,\% on every city. TP-Lit and OD also respond (23 to 69\,\%); VCC is weakly responsive (10 to 13\,\%). Multi-FR is essentially $w_r$-insensitive ($-1\%$ to $-4\%$ on Pittsburgh and Chicago), and degenerates to fleet$=0$ on Manhattan and LA at $w_r{=}10$. This is a known artifact of the calibrated $M = 10{,}000 \cdot w_r$ objective: at high $w_r$ both the per-passenger reward and the wait term scale up, and the master can satisfy load-factor and trip-budget constraints either by selecting $\sim 4$ trips or by selecting zero; on the smaller cities the empty corner sometimes wins. A non-degenerate Multi-FR at $w_r{=}10$ on the smaller cities would require capping $T_{\max}$ at $\sim 1$\,h or relaxing the load-factor floor.

\paragraph{Density scaling ($\Lambda$).} Three regimes emerge as $\Lambda$ rises from 50 to 300 at $w_r{=}1$. (i) Mohring (decreasing user time): SP-Lit drops 20--54\,\%, VCC drops 10--17\,\%, TP-Lit roughly flat. (ii) Anti-Mohring (increasing user time): Joint-Nom/DRO rise 19--34\,\% across all four cities because their Newton $T^*$ \emph{lengthens} with $\Lambda$ to keep tour cost competitive (a CA-literature prediction); OD on Pittsburgh ($+39\%$) and Manhattan ($+144\%$) saturates fleet queueing. (iii) $\Lambda$-insensitive: Multi-FR ($\pm 1\%$ everywhere) because additional demand fills the same already-scheduled trips at the locked $T_{\max}$ headway.

\paragraph{Where each baseline wins.} The cross-city cloud localises the operational choice. Multi-FR is the cheap-provider option (provider $20$--$30$, user time $\sim 1.6$\,h), tolerable only when riders accept long headways (rural or low-demand microtransit). Joint-Nom/DRO and SP-Lit are the provider-balanced option (provider 40--135, user time 0.4--0.9\,h at $w_r{=}1$); the DRO premium over Nom is small ($\sim 0$ to $15\%$) on population-weighted demand. TP-Lit, VCC, and OD are the user-time-optimised option (provider 90--1{,}220, user time 0.07--0.38\,h); cost rises super-linearly with topology irregularity, so the deployment city sharply affects the recommended choice.

\begin{figure}[t]
\centering
\includegraphics[width=\linewidth]{figures/external baselines/designs_four_cities}
\caption{Service designs across the four cities at $\Lambda{=}150$, $w_r{=}1$, with $K{=}3$ fixed for visualisation. Rows are cities (Pittsburgh, Chicago, Manhattan, LA); columns are baselines (Joint-Nom, Joint-DRO, SP-Lit, Multi-FR, TP-Lit, VCC, OD). Route polylines follow road shortest paths. Multi-FR shows selected reference loops; TP-Lit shows three dispatch waves; VCC shows walk halos; OD shows depot plus sample radial trips. Only Joint-Nom and Joint-DRO optimise the depot.}
\label{fig:four-city-designs}
\end{figure}

\paragraph{Implication.} The synthetic-grid conclusions (\S\ref{sec:cloud-analysis}) survive on real cities: the qualitative ranking is topology-robust, and the architectural distinctions remain visible in both the cloud and the design maps. The absolute position of the user-time regime varies with topology by a factor of 2 to 3 across cities, so deployment recommendations should be made on the actual target city, not extrapolated from a synthetic instance. Multi-FR remains the main outlier: its paper-faithful coverage reward yields a competitive provider cost but a long-headway design that pins user time near $T_{\max}/2$ regardless of regime. Reducing $T_{\max}$ or relaxing the load-factor floor would let Multi-FR enter a more user-favorable region of the Pareto plane, at the cost of departing from the paper's calibration.
 where needed.
% Requires: amsmath, booktabs, array.
% Figure paths assume Overleaf layout: figures/external baselines/<name>
% =============================================================================

\section{External Baselines}\label{sec:external-baselines}

We benchmark our service-design framework against seven external baselines drawn from three families: (i) partition with continuous approximation (\textbf{Joint-Nom}, \textbf{Joint-DRO}, \textbf{SP-Lit (Carlsson)}); (ii) routing-based (\textbf{Multi-FR (Jacquillat)}, \textbf{TP-Lit (Liu)}); (iii) operational extremes (\textbf{VCC}, \textbf{OD/fully-on-demand}). Each baseline sweeps its own structural hyperparameters internally per regime point and picks the combination minimising the aggregate cost $C_{\mathrm{prov}} + w_r \cdot \bar{u} \cdot \Lambda$. The comparison is stress-tested at three tiers: \S\ref{sec:cloud-analysis} synthetic $5{\times}5$ grid; \S\ref{sec:real-cloud} a 25-block-group Pittsburgh slice with real road distances and ACS population demand; \S\ref{sec:multi-city} four U.S.\ cities spanning topology archetypes (Pittsburgh, Chicago, Manhattan, LA). To make the regime response readable, each cloud is presented in two faceted views: one figure with $\Lambda$ fixed per panel and $w_r$ varying within the panel (shows how each baseline responds to the user-time weight at a given demand level), and a companion figure with $w_r$ fixed per panel and $\Lambda$ varying within the panel (shows how each baseline responds to demand at a given user-time weight).

\subsection{Analysis on the synthetic instance}\label{sec:cloud-analysis}

\begin{figure}[t]
\centering
\includegraphics[width=0.92\linewidth]{figures/external baselines/cloud_synthetic}
\caption{Performance cloud on the $5{\times}5$ synthetic instance ($\Lambda \in \{50, 150, 300\}$\,pax/h, $w_r \in \{1, 3, 10\}$, $w_v{=}1$, $c_f{=}4.0$\,\$/veh-h). Marker size $\propto \Lambda$, fill shade $\propto w_r$ (lighter $=$ larger $w_r$), shaded hull is each baseline's spread over the $3{\times}3$ regime grid.}
\label{fig:cloud}
\end{figure}

\begin{figure}[t]
\centering
\includegraphics[width=\linewidth]{figures/external baselines/cloud_synthetic_by_lambda}
\caption{Synthetic cloud, fixed $\Lambda$ per panel, $w_r$ varying within panel (fill shade encodes $w_r$, lighter $=$ smaller $w_r$). Lines connect the $w_r$ sweep for each baseline. Baselines whose design is genuinely $w_r$-insensitive at a given $\Lambda$ (TP-Lit at every $\Lambda$, VCC and OD at low $\Lambda$) have all three $w_r$ rows at identical (provider, user) coordinates; markers are then jittered slightly in $\log x$ so all three remain visible.}
\label{fig:cloud-by-lambda}
\end{figure}

\begin{figure}[t]
\centering
\includegraphics[width=\linewidth]{figures/external baselines/cloud_synthetic_by_wr}
\caption{Synthetic cloud, fixed $w_r$ per panel, $\Lambda$ varying within panel (marker size $\propto \Lambda$). Lines connect the $\Lambda$ sweep for each baseline.}
\label{fig:cloud-by-wr}
\end{figure}

Figure~\ref{fig:cloud} plots each baseline's provider cost rate (log scale) against average user time per rider across the full $3 \times 3$ regime grid. Figures~\ref{fig:cloud-by-lambda} and~\ref{fig:cloud-by-wr} decompose the same data by holding one regime parameter fixed per panel. The by-$\Lambda$ view exposes each method's $w_r$ trajectory: partition methods (Joint-Nom, Joint-DRO, SP-Lit) sweep from cheap-provider/long-wait at $w_r{=}1$ to expensive-provider/short-wait at $w_r{=}10$, while TP-Lit, VCC, and OD slide only along the provider axis at roughly fixed user time, and Multi-FR barely moves. The by-$w_r$ view exposes each method's $\Lambda$ trajectory: Multi-FR shows the strongest Mohring effect (user time drops from $\approx 1.07$\,h at $\Lambda{=}50$ to $\approx 0.56$\,h at $\Lambda{=}300$, $w_r{=}1$), partition methods are weakly anti-Mohring (Joint's $T^*$ lengthens with $\Lambda$), and OD/VCC have nearly flat user time (the fleet scales up with $\Lambda$ to absorb demand, so all $\Lambda$ variation appears horizontally on the provider axis).

\subsubsection{Per-baseline sweep grids}

\begin{center}
\small
\renewcommand{\arraystretch}{1.15}
\begin{tabular}{lll}
\toprule
Baseline & Swept hyperparameters & Selection \\
\midrule
Joint-Nom  & $K \in \{2,\dots,6\}$ (random-search partition) & best $K$ by aggregate cost \\
Joint-DRO  & $K \in \{2,\dots,6\}$, $\varepsilon = 0.5$ fixed & best $K$ by aggregate cost \\
SP-Lit     & $K \in \{2,\dots,6\}$ $\times$ $T_{\mathrm{fixed}} \in \{0.1, 0.2, 0.3, 0.5, 0.8, 1.0\}$ & best $(K, T)$ \\
Multi-FR   & (fleet, $\delta$) $\in$ $\{(3,0), (5,0), (8,0), (12,0)\}$, internal $T$-grid & best combo \\
TP-Lit     & $\mathtt{max\_fleet} \in \{20, 50, 100\}$; 25 candidate stops & best fleet \\
VCC        & fleet $\in \{5,10,20,30,50\} \times \delta_w \in \{0.1, 0.3, 0.5\}$\,km & best (fleet, $\delta_w$) \\
OD         & fleet $\in \{5,10,15,20,30,50,80\}$ with queueing & best fleet \\
\bottomrule
\end{tabular}
\end{center}

Joint-DRO holds $\varepsilon{=}0.5$ fixed because it is a robustness knob, not a provider/user Pareto knob.

\subsubsection{Joint-Nom and Joint-DRO}

Joint methods occupy the left band of the cloud (provider cost 28 to 136), spanning a wide user-time range (0.11 to 0.78\,h). Two design knobs respond to the regime: the number of districts $K$ (swept externally over $\{2,\dots,6\}$ by the harness, best $K$ picked per regime point by aggregate cost) and the per-district dispatch interval $T^*$ (computed internally by Newton's method for each candidate $K$).

\begin{center}
\small
\begin{tabular}{cccccc}
\toprule
$\Lambda$ & $w_r$ & $K$ & $T^*$ (h) & User time (h) & Provider cost \\
\midrule
 50 &  1 & 4 & 0.88 & 0.60 &  28 \\
 50 & 10 & 5 & 0.13 & 0.12 &  50 \\
150 &  1 & 6 & 0.99 & 0.70 &  38 \\
150 & 10 & 6 & 0.12 & 0.16 &  95 \\
300 &  1 & 6 & 1.00 & 0.78 &  43 \\
300 & 10 & 6 & 0.13 & 0.19 & 124 \\
\bottomrule
\end{tabular}
\end{center}

$T^*$ is highly $w_r$-responsive: at $w_r{=}1$ the optimiser values rider time cheaply and picks $T^* \approx 0.9$ to $1.0$\,h (heavy batching), while at $w_r{=}10$ it drops to $\approx 0.13$\,h (near real-time dispatch) at 3 to 4$\times$ the provider cost. $K$ increases with both $w_r$ and $\Lambda$ (from $K{=}4$ at the cheap end to $K{=}6$ at the expensive end). User time barely responds to $\Lambda$ at low $w_r$ because the Newton $T^*$ formula uses a per-rider wait term $w_r \cdot T$ rather than aggregate $w_r \cdot T \cdot \Lambda$, and the $K$ sweep partially compensates by adding districts at higher $\Lambda$.

Joint-DRO separates from Joint-Nom only at high $\Lambda$ and $w_r$: at $(\Lambda{=}300, w_r{=}10)$, DRO picks fleet 17 / provider cost 136 versus Nom's fleet 15 / provider cost 124 ($\sim 10\%$ premium). Under uniform demand the DRO premium stays modest; non-uniform demand amplifies it.

\subsubsection{Multi-FR (Jacquillat)}

Multi-FR occupies the mid-left region and is the only method with $\Lambda$-decreasing user time on the synthetic grid (the Mohring effect): more riders justify shorter dispatch headways on the same fixed-route backbone, so $T^*$ scales roughly as $1/\sqrt{\Lambda}$. Fleet 3 wins at every synthetic regime point; fleet 5 is available but rarely preferred.

\begin{center}
\small
\begin{tabular}{cccccc}
\toprule
$\Lambda$ & $w_r$ & $T$ (h) & Wait (h) & User time (h) & Provider cost \\
\midrule
 50 &  1 & 1.75 & 0.83 & 1.04 & 30 \\
150 &  1 & 0.81 & 0.40 & 0.62 & 40 \\
300 &  1 & 0.69 & 0.35 & 0.56 & 44 \\
300 & 10 & 0.38 & 0.19 & 0.39 & 69 \\
\bottomrule
\end{tabular}
\end{center}

Multi-FR is one of the few baselines with non-zero walk time ($\approx 0.07$ to $0.12$\,h on the grid), reflecting the rider-to-checkpoint distance. The evaluator assigns each sampled rider to the nearest realised service point globally rather than to a partition-assigned route.

\subsubsection{SP-Lit (Carlsson)}

SP-Lit sits in the centre-right of the cloud, shifted right of Joint despite the same partition-plus-BHH structure. At every regime point its provider cost is 1.6 to 3.1$\times$ Joint-Nom's. The gap is a partition-quality difference: SP-Lit's deterministic ham-sandwich bisection balances the continuous measures $\mu_1$ (spatial) and $\mu_2$ (radial workload) but does not directly minimise per-district cost; Joint's random search with local improvement finds lower-$\alpha_i$ partitions for the same $K$.

\subsubsection{TP-Lit (Liu)}\label{sec:tplit-cloud}

TP-Lit occupies a narrow horizontal band at user time $\approx 0.20$ to $0.23$\,h with provider cost 233 to 244. On the synthetic instance the design is essentially $w_r$-insensitive: the ADP+VRP method ran at its library defaults ($T{=}0.25$\,h, fleet automatically sized by APT to 22), the three $w_r$ rows at each $\Lambda$ are bit-identical, and the three visible markers per $\Lambda$ panel sit on the same coordinate (small log-$x$ jitter applied for visibility). The real-city sweeps (\S\ref{sec:real-cloud}, \S\ref{sec:multi-city}) do add an external $T_{\mathrm{fixed}} \in \{0.15, 0.5\}$\,h $\times$ $\mathtt{max\_fleet} \in \{30, 100\}$ wrapper sweep and pick the best 4-combo by aggregate cost, but the chosen combo is the same at every $w_r$ at most regime points (e.g.\ all three $w_r$ at $\Lambda{=}150$ Pittsburgh pick $T{=}0.5$, $\mathtt{max\_fleet}{=}30$); a $T{=}0.15$ switch only occurs at $w_r{=}10$ on $\Lambda \in \{50, 300\}$. The reason is structural: TP-Lit's inner objective (ADP delivery-time minimisation) does not contain $w_r$, so only the outer wrapper's aggregate-cost selection responds to $w_r$ at all, and that response is quantised by the 4-combo grid. Maintaining a sizeable fleet for 25 candidate stops makes the fleet-cost term alone ($c_f \times \mathtt{fleet}$) account for $\sim 35\%$ of provider cost on the synthetic.

\subsubsection{VCC and OD}

VCC and OD form the operational-extremes panel (Fig.~\ref{fig:cloud-split}b). Both achieve user time below 0.16\,h on the synthetic grid but at provider cost 1 to 2 orders of magnitude above Joint-Nom. VCC sweeps fleet ($\{5,10,20,30,50\}$) and walk tolerance $\delta_w \in \{0.1,0.3,0.5\}$\,km; the optimiser picks larger fleets at higher $\Lambda$ to keep user time constant ($\approx 0.15$\,h), producing a horizontal trajectory in the cloud. OD serves each request individually (depot $\to$ destination $\to$ depot) and is the cost ceiling on provider expense and the cost floor on user time: every other baseline is interpreted relative to OD's extreme. At low $w_r$, OD tolerates moderate queueing (fleet 5--15); at $w_r{=}10$ it upgrades to 50 vehicles at $\Lambda{=}300$.

\subsubsection{Cross-method synthesis}

No single method dominates across the regime grid. The approximate Pareto frontier runs from Joint-Nom/DRO (cheapest provider, highest user time at low $w_r$) through Multi-FR (Mohring-aided at high $\Lambda$) to VCC and OD (lowest user time, highest provider cost). SP-Lit and TP-Lit are dominated on the frontier: SP-Lit by Joint-Nom (same structure, better partitions) and TP-Lit by VCC (lower user time at comparable provider cost). The clearest architectural distinction is the response to $\Lambda$: Multi-FR's user time \emph{decreases} (Mohring); Joint/SP-Lit barely respond (the $K$ sweep absorbs the demand change); VCC/OD are nearly constant in user time (fleet scales with $\Lambda$); TP-Lit is perfectly flat (its design is $\Lambda$-insensitive in this range).

\subsection{Pittsburgh real-data validation}\label{sec:real-cloud}

\begin{figure}[t]
\centering
\includegraphics[width=\linewidth]{figures/external baselines/cloud_pittsburgh}
\caption{Pittsburgh real-data cloud, 25 census block groups in Allegheny County, road-network shortest paths everywhere (optimisation and evaluation), ACS 2023 population-weighted demand. Same regime grid and same axes as Fig.~\ref{fig:cloud}.}
\label{fig:cloud-real}
\end{figure}

\begin{figure}[t]
\centering
\includegraphics[width=\linewidth]{figures/external baselines/cloud_pittsburgh_by_lambda}
\caption{Pittsburgh cloud, fixed $\Lambda$ per panel, $w_r$ varying.}
\label{fig:cloud-real-by-lambda}
\end{figure}

\begin{figure}[t]
\centering
\includegraphics[width=\linewidth]{figures/external baselines/cloud_pittsburgh_by_wr}
\caption{Pittsburgh cloud, fixed $w_r$ per panel, $\Lambda$ varying.}
\label{fig:cloud-real-by-wr}
\end{figure}

We re-run the same regime grid on 25 census block groups around the centroid of HCT's Pittsburgh service area. Demand is ACS-population-weighted ($p_b \in [0.014, 0.062]$ versus the synthetic $p_b{=}1/25$); distances are OSMnx road-network shortest paths (2{,}931 nodes, 7{,}519 edges) everywhere a baseline supports them (Joint-Nom, Joint-DRO, SP-Lit in optimisation; all baselines in evaluation).

\paragraph{Headline numbers at $\Lambda{=}150$, $w_r{=}1$.}
\begin{center}
\small
\begin{tabular}{lcc}
\toprule
Baseline & Provider cost & User time (h) \\
\midrule
Multi-FR (Jacquillat)  &  29.5 & 1.78 \\
Joint-Nom              &  58.6 & 0.92 \\
Joint-DRO              &  59.5 & 0.91 \\
SP-Lit (Carlsson)      & 134.9 & 0.70 \\
TP-Lit (Liu)           & 216.2 & 0.38 \\
VCC                    & 718.1 & 0.21 \\
OD                     & 1220.2 & 0.23 \\
\bottomrule
\end{tabular}
\end{center}

\paragraph{Findings.} Three patterns from the synthetic analysis carry over and one new one emerges. (i) Provider costs scale up across the board because real road distances are longer than Euclidean (typical road-to-Euclidean factor 1.3 to 1.5 in Pittsburgh's grid); OD's provider cost more than triples, while partition methods absorb most of the increase via their road-aware in-district TSPs. (ii) The Pareto ranking is preserved: same ordering Joint $\to$ SP-Lit $\to$ TP-Lit $\to$ VCC $\to$ OD on provider cost, with the reverse on user time. (iii) Joint's $K$ choice tracks regime monotonically ($K{=}4$ at $\Lambda{=}50, w_r{=}1$, $K{=}6$ at $\Lambda{=}300, w_r{=}10$). (iv) Non-uniform demand sharpens the DRO-vs-Nom premium: at $(\Lambda{=}300, w_r{=}10)$, DRO costs $\sim 25\%$ more than Nom (versus $\sim 10\%$ on uniform synthetic), matching the prediction that DRO's worst-case partition diverges more aggressively when demand is spatially non-uniform.

\begin{figure}[t]
\centering
\includegraphics[width=\linewidth]{figures/external baselines/designs_pittsburgh}
\caption{Pittsburgh service designs at $\Lambda{=}150$, $w_r{=}1$, with $K{=}3$ fixed for visualisation (the cloud sweep picks $K$ per regime). Route polylines follow road shortest paths. Partition methods (cols 1 to 3) show contiguous district fills with exact-TSP depot tours. Multi-FR (col 4) shows selected reference loops. TP-Lit (col 5) shows three independent dispatch waves (blues, reds, greens) with realised VRP routes per vehicle. VCC (col 6) shows walk-radius halos ($\delta_w{=}0.3$\,km). OD (col 7) shows depot plus sample on-demand radial trips.}
\label{fig:designs-map}
\end{figure}

The map overlay confirms that the architectural distinctions are also visible in geometry: Joint-Nom and Joint-DRO produce contiguous random-search partitions with a Joint-Nom-optimised depot; SP-Lit's ham-sandwich partition is more axis-aligned; Multi-FR's selected reference loops radiate from the centroid depot; VCC overlays per-block walk halos; OD shows pure depot-anchored trips.

\subsection{Cross-topology validation across four U.S.\ cities}\label{sec:multi-city}

\begin{figure}[t]
\centering
\includegraphics[width=\linewidth]{figures/external baselines/cloud_four_cities}
\caption{Cross-city baseline performance clouds on the shared regime grid. Each panel is one city; markers and shading as in Fig.~\ref{fig:cloud}.}
\label{fig:four-city-cloud}
\end{figure}

\begin{figure}[t]
\centering
\includegraphics[width=\linewidth]{figures/external baselines/cloud_four_cities_by_lambda}
\caption{Cross-city clouds, fixed $\Lambda$ per column, $w_r$ varying within panel. Rows are cities.}
\label{fig:four-city-cloud-by-lambda}
\end{figure}

\begin{figure}[t]
\centering
\includegraphics[width=\linewidth]{figures/external baselines/cloud_four_cities_by_wr}
\caption{Cross-city clouds, fixed $w_r$ per column, $\Lambda$ varying within panel. Rows are cities.}
\label{fig:four-city-cloud-by-wr}
\end{figure}

We re-run the regime sweep on four cities chosen for road-topology contrast: Pittsburgh (irregular hilly, river-cut), Chicago Loop (strict cardinal grid), Manhattan Midtown (avenue-and-street grid), and Los Angeles DTLA (freeway-fragmented grid). Each city uses real TIGER 2023 block groups, ACS-population demand, and full road-network shortest paths in both optimisation and evaluation.

\paragraph{Service-region geometry.} BG count is held at 25 (50 for Manhattan, see below); TIGER block groups are population-sized, so denser cities pack the same count into smaller polygons.

\begin{center}
\small
\begin{tabular}{lcrrrr}
\toprule
City & BGs & Area (km$^2$) & BBox W$\times$H (km) & Diag (km) & Density (pop/km$^2$) \\
\midrule
Pittsburgh & 25 & 24.5 & $7.4 \times 7.2$ & 10.3 &  1{,}003 \\
Chicago    & 25 &  4.6 & $2.3 \times 3.5$ &  4.1 & 13{,}937 \\
Manhattan  & 50 &  2.4 & $2.0 \times 2.5$ &  3.2 & 18{,}891 \\
LA         & 25 &  4.3 & $2.3 \times 2.9$ &  3.7 &  8{,}755 \\
\bottomrule
\end{tabular}
\end{center}

Pittsburgh's 25 BGs cover roughly $10{\times}$ more area than Manhattan's 50, so cross-city cost comparisons reflect both topology and area. Manhattan is enlarged from the default 25 BGs ($1.1$\,km$^2$, too small for FR-Detour's candidate-line library to produce feasible trips) to 50 BGs ($2.4$\,km$^2$); the other three cities are kept at 25 BGs.

\paragraph{Headline numbers at $\Lambda{=}150$, $w_r{=}1$.}
\begin{center}
\small
\begin{tabular}{lrrrr}
\toprule
Baseline & Pittsburgh & Chicago & Manhattan & LA \\
\midrule
\multicolumn{5}{l}{\emph{Provider cost}} \\
Joint-Nom              &   58.6 &   47.6 &   47.9 &   35.5 \\
Joint-DRO              &   59.5 &   45.8 &   54.6 &   36.7 \\
SP-Lit (Carlsson)      &  134.9 &   81.1 &   73.7 &   58.5 \\
Multi-FR (Jacquillat)  &   29.5 &   21.4 &   20.5 &   20.5 \\
TP-Lit (Liu)           &  216.2 &  103.1 &  104.7 &   91.5 \\
VCC                    &  718.1 &  241.9 &  199.7 &  193.3 \\
OD                     & 1220.2 &  476.4 &  417.5 &  343.8 \\
\midrule
\multicolumn{5}{l}{\emph{User time (h per rider)}} \\
Joint-Nom              & 0.92 & 0.64 & 0.58 & 0.56 \\
Joint-DRO              & 0.91 & 0.68 & 0.60 & 0.56 \\
SP-Lit (Carlsson)      & 0.70 & 0.41 & 0.39 & 0.41 \\
Multi-FR (Jacquillat)  & 1.78 & 1.61 & 1.58 & 1.59 \\
TP-Lit (Liu)           & 0.38 & 0.31 & 0.31 & 0.30 \\
VCC                    & 0.21 & 0.12 & 0.11 & 0.09 \\
OD                     & 0.23 & 0.07 & 0.18 & 0.08 \\
\bottomrule
\end{tabular}
\end{center}

\paragraph{Multi-FR calibration.} Multi-FR uses the paper-faithful coverage reward $M = 10{,}000 \cdot w_r$ (Jacquillat et al.\ Table EC.4). With this calibration the master picks fleet $\in \{4, 6, 9, 12\}$ across the regime grid (versus the degenerate fleet $=1$ under our earlier $M = 10 \cdot w_r$), and walk time drops from 0.4 to 0.5\,h down to $\sim 0.07$\,h. However, the master still selects the longest scheduled headway in our $T$-grid ($T_{\max}{=}3$\,h, realised wait $\approx 1.5$\,h) at every regime point because $M{=}10^4 \cdot w_r$ dwarfs the per-passenger wait term $w_r \cdot T/2 \cdot \mathrm{demand}$ by $\sim 6{,}700{\times}$. Multi-FR therefore lands at a many-trips, full-coverage, long-headway design point on real cities. At $w_r{=}10$ on Manhattan and LA the master sometimes flips to an empty-solution corner of the feasible set; see the discussion below.

\paragraph{Cross-method ranking and Pareto positions.} The provider-cost ordering Multi-FR $<$ Joint $<$ SP-Lit $<$ TP-Lit $<$ VCC $<$ OD holds at every $(\Lambda, w_r)$ point on all four cities; the user-time ordering is essentially the reverse. Multi-FR consistently occupies the cheapest-provider corner (provider 20--30, user time 1.6--1.8\,h); OD the opposite (provider 344--1{,}220, user time 0.07--0.23\,h). The Pareto frontier is dense: at $(\Lambda{=}150, w_r{=}1)$, five to seven of the seven baselines are non-dominated on each city. The architectural distinctions between baselines correspond to genuinely different operating points on the same Pareto plane.

\paragraph{Topology controls absolute cost.} Pittsburgh (25 BGs, 24.5\,km$^2$, irregular) inflates absolute provider cost relative to Manhattan (50 BGs, 2.4\,km$^2$, dense grid): Joint-Nom $59$ vs $48$ ($1.2{\times}$), SP-Lit $135$ vs $74$ ($1.8{\times}$), TP-Lit $216$ vs $105$ ($2.1{\times}$), OD $1{,}220$ vs $418$ ($2.9{\times}$). Methods that pay distance per request (OD, VCC) amplify the geometric handicap; partition methods absorb most of it via a shared in-district TSP. Multi-FR is essentially topology-invariant in provider cost (clustering at 20--30 across all four cities) because its long-headway design is dispatch-frequency-bound, not distance-bound. The OD-to-Joint multiplier rises with topology complexity: $9{\times}$ on Manhattan, 10--13$\times$ on Chicago and LA, $21{\times}$ on Pittsburgh.

\paragraph{Regime responsiveness ($w_r$).} As $w_r$ rises from 1 to 10 at $\Lambda{=}150$, partition methods (Joint-Nom/DRO, SP-Lit) drop user time by 66 to 77\,\% on every city. TP-Lit and OD also respond (23 to 69\,\%); VCC is weakly responsive (10 to 13\,\%). Multi-FR is essentially $w_r$-insensitive ($-1\%$ to $-4\%$ on Pittsburgh and Chicago), and degenerates to fleet$=0$ on Manhattan and LA at $w_r{=}10$. This is a known artifact of the calibrated $M = 10{,}000 \cdot w_r$ objective: at high $w_r$ both the per-passenger reward and the wait term scale up, and the master can satisfy load-factor and trip-budget constraints either by selecting $\sim 4$ trips or by selecting zero; on the smaller cities the empty corner sometimes wins. A non-degenerate Multi-FR at $w_r{=}10$ on the smaller cities would require capping $T_{\max}$ at $\sim 1$\,h or relaxing the load-factor floor.

\paragraph{Density scaling ($\Lambda$).} Three regimes emerge as $\Lambda$ rises from 50 to 300 at $w_r{=}1$. (i) Mohring (decreasing user time): SP-Lit drops 20--54\,\%, VCC drops 10--17\,\%, TP-Lit roughly flat. (ii) Anti-Mohring (increasing user time): Joint-Nom/DRO rise 19--34\,\% across all four cities because their Newton $T^*$ \emph{lengthens} with $\Lambda$ to keep tour cost competitive (a CA-literature prediction); OD on Pittsburgh ($+39\%$) and Manhattan ($+144\%$) saturates fleet queueing. (iii) $\Lambda$-insensitive: Multi-FR ($\pm 1\%$ everywhere) because additional demand fills the same already-scheduled trips at the locked $T_{\max}$ headway.

\paragraph{Where each baseline wins.} The cross-city cloud localises the operational choice. Multi-FR is the cheap-provider option (provider $20$--$30$, user time $\sim 1.6$\,h), tolerable only when riders accept long headways (rural or low-demand microtransit). Joint-Nom/DRO and SP-Lit are the provider-balanced option (provider 40--135, user time 0.4--0.9\,h at $w_r{=}1$); the DRO premium over Nom is small ($\sim 0$ to $15\%$) on population-weighted demand. TP-Lit, VCC, and OD are the user-time-optimised option (provider 90--1{,}220, user time 0.07--0.38\,h); cost rises super-linearly with topology irregularity, so the deployment city sharply affects the recommended choice.

\begin{figure}[t]
\centering
\includegraphics[width=\linewidth]{figures/external baselines/designs_four_cities}
\caption{Service designs across the four cities at $\Lambda{=}150$, $w_r{=}1$, with $K{=}3$ fixed for visualisation. Rows are cities (Pittsburgh, Chicago, Manhattan, LA); columns are baselines (Joint-Nom, Joint-DRO, SP-Lit, Multi-FR, TP-Lit, VCC, OD). Route polylines follow road shortest paths. Multi-FR shows selected reference loops; TP-Lit shows three dispatch waves; VCC shows walk halos; OD shows depot plus sample radial trips. Only Joint-Nom and Joint-DRO optimise the depot.}
\label{fig:four-city-designs}
\end{figure}

\paragraph{Implication.} The synthetic-grid conclusions (\S\ref{sec:cloud-analysis}) survive on real cities: the qualitative ranking is topology-robust, and the architectural distinctions remain visible in both the cloud and the design maps. The absolute position of the user-time regime varies with topology by a factor of 2 to 3 across cities, so deployment recommendations should be made on the actual target city, not extrapolated from a synthetic instance. Multi-FR remains the main outlier: its paper-faithful coverage reward yields a competitive provider cost but a long-headway design that pins user time near $T_{\max}/2$ regardless of regime. Reducing $T_{\max}$ or relaxing the load-factor floor would let Multi-FR enter a more user-favorable region of the Pareto plane, at the cost of departing from the paper's calibration.
 where needed.
% Requires: amsmath, booktabs, array.
% Figure paths assume Overleaf layout: figures/external baselines/<name>
% =============================================================================

\section{External Baselines}\label{sec:external-baselines}

We benchmark our service-design framework against seven external baselines drawn from three families: (i) partition with continuous approximation (\textbf{Joint-Nom}, \textbf{Joint-DRO}, \textbf{SP-Lit (Carlsson)}); (ii) routing-based (\textbf{Multi-FR (Jacquillat)}, \textbf{TP-Lit (Liu)}); (iii) operational extremes (\textbf{VCC}, \textbf{OD/fully-on-demand}). Each baseline sweeps its own structural hyperparameters internally per regime point and picks the combination minimising the aggregate cost $C_{\mathrm{prov}} + w_r \cdot \bar{u} \cdot \Lambda$. The comparison is stress-tested at three tiers: \S\ref{sec:cloud-analysis} synthetic $5{\times}5$ grid; \S\ref{sec:real-cloud} a 25-block-group Pittsburgh slice with real road distances and ACS population demand; \S\ref{sec:multi-city} four U.S.\ cities spanning topology archetypes (Pittsburgh, Chicago, Manhattan, LA). To make the regime response readable, each cloud is presented in two faceted views: one figure with $\Lambda$ fixed per panel and $w_r$ varying within the panel (shows how each baseline responds to the user-time weight at a given demand level), and a companion figure with $w_r$ fixed per panel and $\Lambda$ varying within the panel (shows how each baseline responds to demand at a given user-time weight).

\subsection{Analysis on the synthetic instance}\label{sec:cloud-analysis}

\begin{figure}[t]
\centering
\includegraphics[width=0.92\linewidth]{figures/external baselines/cloud_synthetic}
\caption{Performance cloud on the $5{\times}5$ synthetic instance ($\Lambda \in \{50, 150, 300\}$\,pax/h, $w_r \in \{1, 3, 10\}$, $w_v{=}1$, $c_f{=}4.0$\,\$/veh-h). Marker size $\propto \Lambda$, fill shade $\propto w_r$ (lighter $=$ larger $w_r$), shaded hull is each baseline's spread over the $3{\times}3$ regime grid.}
\label{fig:cloud}
\end{figure}

\begin{figure}[t]
\centering
\includegraphics[width=\linewidth]{figures/external baselines/cloud_synthetic_by_lambda}
\caption{Synthetic cloud, fixed $\Lambda$ per panel, $w_r$ varying within panel (fill shade encodes $w_r$, lighter $=$ smaller $w_r$). Lines connect the $w_r$ sweep for each baseline. Baselines whose design is genuinely $w_r$-insensitive at a given $\Lambda$ (TP-Lit at every $\Lambda$, VCC and OD at low $\Lambda$) have all three $w_r$ rows at identical (provider, user) coordinates; markers are then jittered slightly in $\log x$ so all three remain visible.}
\label{fig:cloud-by-lambda}
\end{figure}

\begin{figure}[t]
\centering
\includegraphics[width=\linewidth]{figures/external baselines/cloud_synthetic_by_wr}
\caption{Synthetic cloud, fixed $w_r$ per panel, $\Lambda$ varying within panel (marker size $\propto \Lambda$). Lines connect the $\Lambda$ sweep for each baseline.}
\label{fig:cloud-by-wr}
\end{figure}

Figure~\ref{fig:cloud} plots each baseline's provider cost rate (log scale) against average user time per rider across the full $3 \times 3$ regime grid. Figures~\ref{fig:cloud-by-lambda} and~\ref{fig:cloud-by-wr} decompose the same data by holding one regime parameter fixed per panel. The by-$\Lambda$ view exposes each method's $w_r$ trajectory: partition methods (Joint-Nom, Joint-DRO, SP-Lit) sweep from cheap-provider/long-wait at $w_r{=}1$ to expensive-provider/short-wait at $w_r{=}10$, while TP-Lit, VCC, and OD slide only along the provider axis at roughly fixed user time, and Multi-FR barely moves. The by-$w_r$ view exposes each method's $\Lambda$ trajectory: Multi-FR shows the strongest Mohring effect (user time drops from $\approx 1.07$\,h at $\Lambda{=}50$ to $\approx 0.56$\,h at $\Lambda{=}300$, $w_r{=}1$), partition methods are weakly anti-Mohring (Joint's $T^*$ lengthens with $\Lambda$), and OD/VCC have nearly flat user time (the fleet scales up with $\Lambda$ to absorb demand, so all $\Lambda$ variation appears horizontally on the provider axis).

\subsubsection{Per-baseline sweep grids}

\begin{center}
\small
\renewcommand{\arraystretch}{1.15}
\begin{tabular}{lll}
\toprule
Baseline & Swept hyperparameters & Selection \\
\midrule
Joint-Nom  & $K \in \{2,\dots,6\}$ (random-search partition) & best $K$ by aggregate cost \\
Joint-DRO  & $K \in \{2,\dots,6\}$, $\varepsilon = 0.5$ fixed & best $K$ by aggregate cost \\
SP-Lit     & $K \in \{2,\dots,6\}$ $\times$ $T_{\mathrm{fixed}} \in \{0.1, 0.2, 0.3, 0.5, 0.8, 1.0\}$ & best $(K, T)$ \\
Multi-FR   & (fleet, $\delta$) $\in$ $\{(3,0), (5,0), (8,0), (12,0)\}$, internal $T$-grid & best combo \\
TP-Lit     & $\mathtt{max\_fleet} \in \{20, 50, 100\}$; 25 candidate stops & best fleet \\
VCC        & fleet $\in \{5,10,20,30,50\} \times \delta_w \in \{0.1, 0.3, 0.5\}$\,km & best (fleet, $\delta_w$) \\
OD         & fleet $\in \{5,10,15,20,30,50,80\}$ with queueing & best fleet \\
\bottomrule
\end{tabular}
\end{center}

Joint-DRO holds $\varepsilon{=}0.5$ fixed because it is a robustness knob, not a provider/user Pareto knob.

\subsubsection{Joint-Nom and Joint-DRO}

Joint methods occupy the left band of the cloud (provider cost 28 to 136), spanning a wide user-time range (0.11 to 0.78\,h). Two design knobs respond to the regime: the number of districts $K$ (swept externally over $\{2,\dots,6\}$ by the harness, best $K$ picked per regime point by aggregate cost) and the per-district dispatch interval $T^*$ (computed internally by Newton's method for each candidate $K$).

\begin{center}
\small
\begin{tabular}{cccccc}
\toprule
$\Lambda$ & $w_r$ & $K$ & $T^*$ (h) & User time (h) & Provider cost \\
\midrule
 50 &  1 & 4 & 0.88 & 0.60 &  28 \\
 50 & 10 & 5 & 0.13 & 0.12 &  50 \\
150 &  1 & 6 & 0.99 & 0.70 &  38 \\
150 & 10 & 6 & 0.12 & 0.16 &  95 \\
300 &  1 & 6 & 1.00 & 0.78 &  43 \\
300 & 10 & 6 & 0.13 & 0.19 & 124 \\
\bottomrule
\end{tabular}
\end{center}

$T^*$ is highly $w_r$-responsive: at $w_r{=}1$ the optimiser values rider time cheaply and picks $T^* \approx 0.9$ to $1.0$\,h (heavy batching), while at $w_r{=}10$ it drops to $\approx 0.13$\,h (near real-time dispatch) at 3 to 4$\times$ the provider cost. $K$ increases with both $w_r$ and $\Lambda$ (from $K{=}4$ at the cheap end to $K{=}6$ at the expensive end). User time barely responds to $\Lambda$ at low $w_r$ because the Newton $T^*$ formula uses a per-rider wait term $w_r \cdot T$ rather than aggregate $w_r \cdot T \cdot \Lambda$, and the $K$ sweep partially compensates by adding districts at higher $\Lambda$.

Joint-DRO separates from Joint-Nom only at high $\Lambda$ and $w_r$: at $(\Lambda{=}300, w_r{=}10)$, DRO picks fleet 17 / provider cost 136 versus Nom's fleet 15 / provider cost 124 ($\sim 10\%$ premium). Under uniform demand the DRO premium stays modest; non-uniform demand amplifies it.

\subsubsection{Multi-FR (Jacquillat)}

Multi-FR occupies the mid-left region and is the only method with $\Lambda$-decreasing user time on the synthetic grid (the Mohring effect): more riders justify shorter dispatch headways on the same fixed-route backbone, so $T^*$ scales roughly as $1/\sqrt{\Lambda}$. Fleet 3 wins at every synthetic regime point; fleet 5 is available but rarely preferred.

\begin{center}
\small
\begin{tabular}{cccccc}
\toprule
$\Lambda$ & $w_r$ & $T$ (h) & Wait (h) & User time (h) & Provider cost \\
\midrule
 50 &  1 & 1.75 & 0.83 & 1.04 & 30 \\
150 &  1 & 0.81 & 0.40 & 0.62 & 40 \\
300 &  1 & 0.69 & 0.35 & 0.56 & 44 \\
300 & 10 & 0.38 & 0.19 & 0.39 & 69 \\
\bottomrule
\end{tabular}
\end{center}

Multi-FR is one of the few baselines with non-zero walk time ($\approx 0.07$ to $0.12$\,h on the grid), reflecting the rider-to-checkpoint distance. The evaluator assigns each sampled rider to the nearest realised service point globally rather than to a partition-assigned route.

\subsubsection{SP-Lit (Carlsson)}

SP-Lit sits in the centre-right of the cloud, shifted right of Joint despite the same partition-plus-BHH structure. At every regime point its provider cost is 1.6 to 3.1$\times$ Joint-Nom's. The gap is a partition-quality difference: SP-Lit's deterministic ham-sandwich bisection balances the continuous measures $\mu_1$ (spatial) and $\mu_2$ (radial workload) but does not directly minimise per-district cost; Joint's random search with local improvement finds lower-$\alpha_i$ partitions for the same $K$.

\subsubsection{TP-Lit (Liu)}\label{sec:tplit-cloud}

TP-Lit occupies a narrow horizontal band at user time $\approx 0.20$ to $0.23$\,h with provider cost 233 to 244. On the synthetic instance the design is essentially $w_r$-insensitive: the ADP+VRP method ran at its library defaults ($T{=}0.25$\,h, fleet automatically sized by APT to 22), the three $w_r$ rows at each $\Lambda$ are bit-identical, and the three visible markers per $\Lambda$ panel sit on the same coordinate (small log-$x$ jitter applied for visibility). The real-city sweeps (\S\ref{sec:real-cloud}, \S\ref{sec:multi-city}) do add an external $T_{\mathrm{fixed}} \in \{0.15, 0.5\}$\,h $\times$ $\mathtt{max\_fleet} \in \{30, 100\}$ wrapper sweep and pick the best 4-combo by aggregate cost, but the chosen combo is the same at every $w_r$ at most regime points (e.g.\ all three $w_r$ at $\Lambda{=}150$ Pittsburgh pick $T{=}0.5$, $\mathtt{max\_fleet}{=}30$); a $T{=}0.15$ switch only occurs at $w_r{=}10$ on $\Lambda \in \{50, 300\}$. The reason is structural: TP-Lit's inner objective (ADP delivery-time minimisation) does not contain $w_r$, so only the outer wrapper's aggregate-cost selection responds to $w_r$ at all, and that response is quantised by the 4-combo grid. Maintaining a sizeable fleet for 25 candidate stops makes the fleet-cost term alone ($c_f \times \mathtt{fleet}$) account for $\sim 35\%$ of provider cost on the synthetic.

\subsubsection{VCC and OD}

VCC and OD form the operational-extremes panel (Fig.~\ref{fig:cloud-split}b). Both achieve user time below 0.16\,h on the synthetic grid but at provider cost 1 to 2 orders of magnitude above Joint-Nom. VCC sweeps fleet ($\{5,10,20,30,50\}$) and walk tolerance $\delta_w \in \{0.1,0.3,0.5\}$\,km; the optimiser picks larger fleets at higher $\Lambda$ to keep user time constant ($\approx 0.15$\,h), producing a horizontal trajectory in the cloud. OD serves each request individually (depot $\to$ destination $\to$ depot) and is the cost ceiling on provider expense and the cost floor on user time: every other baseline is interpreted relative to OD's extreme. At low $w_r$, OD tolerates moderate queueing (fleet 5--15); at $w_r{=}10$ it upgrades to 50 vehicles at $\Lambda{=}300$.

\subsubsection{Cross-method synthesis}

No single method dominates across the regime grid. The approximate Pareto frontier runs from Joint-Nom/DRO (cheapest provider, highest user time at low $w_r$) through Multi-FR (Mohring-aided at high $\Lambda$) to VCC and OD (lowest user time, highest provider cost). SP-Lit and TP-Lit are dominated on the frontier: SP-Lit by Joint-Nom (same structure, better partitions) and TP-Lit by VCC (lower user time at comparable provider cost). The clearest architectural distinction is the response to $\Lambda$: Multi-FR's user time \emph{decreases} (Mohring); Joint/SP-Lit barely respond (the $K$ sweep absorbs the demand change); VCC/OD are nearly constant in user time (fleet scales with $\Lambda$); TP-Lit is perfectly flat (its design is $\Lambda$-insensitive in this range).

\subsection{Pittsburgh real-data validation}\label{sec:real-cloud}

\begin{figure}[t]
\centering
\includegraphics[width=\linewidth]{figures/external baselines/cloud_pittsburgh}
\caption{Pittsburgh real-data cloud, 25 census block groups in Allegheny County, road-network shortest paths everywhere (optimisation and evaluation), ACS 2023 population-weighted demand. Same regime grid and same axes as Fig.~\ref{fig:cloud}.}
\label{fig:cloud-real}
\end{figure}

\begin{figure}[t]
\centering
\includegraphics[width=\linewidth]{figures/external baselines/cloud_pittsburgh_by_lambda}
\caption{Pittsburgh cloud, fixed $\Lambda$ per panel, $w_r$ varying.}
\label{fig:cloud-real-by-lambda}
\end{figure}

\begin{figure}[t]
\centering
\includegraphics[width=\linewidth]{figures/external baselines/cloud_pittsburgh_by_wr}
\caption{Pittsburgh cloud, fixed $w_r$ per panel, $\Lambda$ varying.}
\label{fig:cloud-real-by-wr}
\end{figure}

We re-run the same regime grid on 25 census block groups around the centroid of HCT's Pittsburgh service area. Demand is ACS-population-weighted ($p_b \in [0.014, 0.062]$ versus the synthetic $p_b{=}1/25$); distances are OSMnx road-network shortest paths (2{,}931 nodes, 7{,}519 edges) everywhere a baseline supports them (Joint-Nom, Joint-DRO, SP-Lit in optimisation; all baselines in evaluation).

\paragraph{Headline numbers at $\Lambda{=}150$, $w_r{=}1$.}
\begin{center}
\small
\begin{tabular}{lcc}
\toprule
Baseline & Provider cost & User time (h) \\
\midrule
Multi-FR (Jacquillat)  &  29.5 & 1.78 \\
Joint-Nom              &  58.6 & 0.92 \\
Joint-DRO              &  59.5 & 0.91 \\
SP-Lit (Carlsson)      & 134.9 & 0.70 \\
TP-Lit (Liu)           & 216.2 & 0.38 \\
VCC                    & 718.1 & 0.21 \\
OD                     & 1220.2 & 0.23 \\
\bottomrule
\end{tabular}
\end{center}

\paragraph{Findings.} Three patterns from the synthetic analysis carry over and one new one emerges. (i) Provider costs scale up across the board because real road distances are longer than Euclidean (typical road-to-Euclidean factor 1.3 to 1.5 in Pittsburgh's grid); OD's provider cost more than triples, while partition methods absorb most of the increase via their road-aware in-district TSPs. (ii) The Pareto ranking is preserved: same ordering Joint $\to$ SP-Lit $\to$ TP-Lit $\to$ VCC $\to$ OD on provider cost, with the reverse on user time. (iii) Joint's $K$ choice tracks regime monotonically ($K{=}4$ at $\Lambda{=}50, w_r{=}1$, $K{=}6$ at $\Lambda{=}300, w_r{=}10$). (iv) Non-uniform demand sharpens the DRO-vs-Nom premium: at $(\Lambda{=}300, w_r{=}10)$, DRO costs $\sim 25\%$ more than Nom (versus $\sim 10\%$ on uniform synthetic), matching the prediction that DRO's worst-case partition diverges more aggressively when demand is spatially non-uniform.

\begin{figure}[t]
\centering
\includegraphics[width=\linewidth]{figures/external baselines/designs_pittsburgh}
\caption{Pittsburgh service designs at $\Lambda{=}150$, $w_r{=}1$, with $K{=}3$ fixed for visualisation (the cloud sweep picks $K$ per regime). Route polylines follow road shortest paths. Partition methods (cols 1 to 3) show contiguous district fills with exact-TSP depot tours. Multi-FR (col 4) shows selected reference loops. TP-Lit (col 5) shows three independent dispatch waves (blues, reds, greens) with realised VRP routes per vehicle. VCC (col 6) shows walk-radius halos ($\delta_w{=}0.3$\,km). OD (col 7) shows depot plus sample on-demand radial trips.}
\label{fig:designs-map}
\end{figure}

The map overlay confirms that the architectural distinctions are also visible in geometry: Joint-Nom and Joint-DRO produce contiguous random-search partitions with a Joint-Nom-optimised depot; SP-Lit's ham-sandwich partition is more axis-aligned; Multi-FR's selected reference loops radiate from the centroid depot; VCC overlays per-block walk halos; OD shows pure depot-anchored trips.

\subsection{Cross-topology validation across four U.S.\ cities}\label{sec:multi-city}

\begin{figure}[t]
\centering
\includegraphics[width=\linewidth]{figures/external baselines/cloud_four_cities}
\caption{Cross-city baseline performance clouds on the shared regime grid. Each panel is one city; markers and shading as in Fig.~\ref{fig:cloud}.}
\label{fig:four-city-cloud}
\end{figure}

\begin{figure}[t]
\centering
\includegraphics[width=\linewidth]{figures/external baselines/cloud_four_cities_by_lambda}
\caption{Cross-city clouds, fixed $\Lambda$ per column, $w_r$ varying within panel. Rows are cities.}
\label{fig:four-city-cloud-by-lambda}
\end{figure}

\begin{figure}[t]
\centering
\includegraphics[width=\linewidth]{figures/external baselines/cloud_four_cities_by_wr}
\caption{Cross-city clouds, fixed $w_r$ per column, $\Lambda$ varying within panel. Rows are cities.}
\label{fig:four-city-cloud-by-wr}
\end{figure}

We re-run the regime sweep on four cities chosen for road-topology contrast: Pittsburgh (irregular hilly, river-cut), Chicago Loop (strict cardinal grid), Manhattan Midtown (avenue-and-street grid), and Los Angeles DTLA (freeway-fragmented grid). Each city uses real TIGER 2023 block groups, ACS-population demand, and full road-network shortest paths in both optimisation and evaluation.

\paragraph{Service-region geometry.} BG count is held at 25 (50 for Manhattan, see below); TIGER block groups are population-sized, so denser cities pack the same count into smaller polygons.

\begin{center}
\small
\begin{tabular}{lcrrrr}
\toprule
City & BGs & Area (km$^2$) & BBox W$\times$H (km) & Diag (km) & Density (pop/km$^2$) \\
\midrule
Pittsburgh & 25 & 24.5 & $7.4 \times 7.2$ & 10.3 &  1{,}003 \\
Chicago    & 25 &  4.6 & $2.3 \times 3.5$ &  4.1 & 13{,}937 \\
Manhattan  & 50 &  2.4 & $2.0 \times 2.5$ &  3.2 & 18{,}891 \\
LA         & 25 &  4.3 & $2.3 \times 2.9$ &  3.7 &  8{,}755 \\
\bottomrule
\end{tabular}
\end{center}

Pittsburgh's 25 BGs cover roughly $10{\times}$ more area than Manhattan's 50, so cross-city cost comparisons reflect both topology and area. Manhattan is enlarged from the default 25 BGs ($1.1$\,km$^2$, too small for FR-Detour's candidate-line library to produce feasible trips) to 50 BGs ($2.4$\,km$^2$); the other three cities are kept at 25 BGs.

\paragraph{Headline numbers at $\Lambda{=}150$, $w_r{=}1$.}
\begin{center}
\small
\begin{tabular}{lrrrr}
\toprule
Baseline & Pittsburgh & Chicago & Manhattan & LA \\
\midrule
\multicolumn{5}{l}{\emph{Provider cost}} \\
Joint-Nom              &   58.6 &   47.6 &   47.9 &   35.5 \\
Joint-DRO              &   59.5 &   45.8 &   54.6 &   36.7 \\
SP-Lit (Carlsson)      &  134.9 &   81.1 &   73.7 &   58.5 \\
Multi-FR (Jacquillat)  &   29.5 &   21.4 &   20.5 &   20.5 \\
TP-Lit (Liu)           &  216.2 &  103.1 &  104.7 &   91.5 \\
VCC                    &  718.1 &  241.9 &  199.7 &  193.3 \\
OD                     & 1220.2 &  476.4 &  417.5 &  343.8 \\
\midrule
\multicolumn{5}{l}{\emph{User time (h per rider)}} \\
Joint-Nom              & 0.92 & 0.64 & 0.58 & 0.56 \\
Joint-DRO              & 0.91 & 0.68 & 0.60 & 0.56 \\
SP-Lit (Carlsson)      & 0.70 & 0.41 & 0.39 & 0.41 \\
Multi-FR (Jacquillat)  & 1.78 & 1.61 & 1.58 & 1.59 \\
TP-Lit (Liu)           & 0.38 & 0.31 & 0.31 & 0.30 \\
VCC                    & 0.21 & 0.12 & 0.11 & 0.09 \\
OD                     & 0.23 & 0.07 & 0.18 & 0.08 \\
\bottomrule
\end{tabular}
\end{center}

\paragraph{Multi-FR calibration.} Multi-FR uses the paper-faithful coverage reward $M = 10{,}000 \cdot w_r$ (Jacquillat et al.\ Table EC.4). With this calibration the master picks fleet $\in \{4, 6, 9, 12\}$ across the regime grid (versus the degenerate fleet $=1$ under our earlier $M = 10 \cdot w_r$), and walk time drops from 0.4 to 0.5\,h down to $\sim 0.07$\,h. However, the master still selects the longest scheduled headway in our $T$-grid ($T_{\max}{=}3$\,h, realised wait $\approx 1.5$\,h) at every regime point because $M{=}10^4 \cdot w_r$ dwarfs the per-passenger wait term $w_r \cdot T/2 \cdot \mathrm{demand}$ by $\sim 6{,}700{\times}$. Multi-FR therefore lands at a many-trips, full-coverage, long-headway design point on real cities. At $w_r{=}10$ on Manhattan and LA the master sometimes flips to an empty-solution corner of the feasible set; see the discussion below.

\paragraph{Cross-method ranking and Pareto positions.} The provider-cost ordering Multi-FR $<$ Joint $<$ SP-Lit $<$ TP-Lit $<$ VCC $<$ OD holds at every $(\Lambda, w_r)$ point on all four cities; the user-time ordering is essentially the reverse. Multi-FR consistently occupies the cheapest-provider corner (provider 20--30, user time 1.6--1.8\,h); OD the opposite (provider 344--1{,}220, user time 0.07--0.23\,h). The Pareto frontier is dense: at $(\Lambda{=}150, w_r{=}1)$, five to seven of the seven baselines are non-dominated on each city. The architectural distinctions between baselines correspond to genuinely different operating points on the same Pareto plane.

\paragraph{Topology controls absolute cost.} Pittsburgh (25 BGs, 24.5\,km$^2$, irregular) inflates absolute provider cost relative to Manhattan (50 BGs, 2.4\,km$^2$, dense grid): Joint-Nom $59$ vs $48$ ($1.2{\times}$), SP-Lit $135$ vs $74$ ($1.8{\times}$), TP-Lit $216$ vs $105$ ($2.1{\times}$), OD $1{,}220$ vs $418$ ($2.9{\times}$). Methods that pay distance per request (OD, VCC) amplify the geometric handicap; partition methods absorb most of it via a shared in-district TSP. Multi-FR is essentially topology-invariant in provider cost (clustering at 20--30 across all four cities) because its long-headway design is dispatch-frequency-bound, not distance-bound. The OD-to-Joint multiplier rises with topology complexity: $9{\times}$ on Manhattan, 10--13$\times$ on Chicago and LA, $21{\times}$ on Pittsburgh.

\paragraph{Regime responsiveness ($w_r$).} As $w_r$ rises from 1 to 10 at $\Lambda{=}150$, partition methods (Joint-Nom/DRO, SP-Lit) drop user time by 66 to 77\,\% on every city. TP-Lit and OD also respond (23 to 69\,\%); VCC is weakly responsive (10 to 13\,\%). Multi-FR is essentially $w_r$-insensitive ($-1\%$ to $-4\%$ on Pittsburgh and Chicago), and degenerates to fleet$=0$ on Manhattan and LA at $w_r{=}10$. This is a known artifact of the calibrated $M = 10{,}000 \cdot w_r$ objective: at high $w_r$ both the per-passenger reward and the wait term scale up, and the master can satisfy load-factor and trip-budget constraints either by selecting $\sim 4$ trips or by selecting zero; on the smaller cities the empty corner sometimes wins. A non-degenerate Multi-FR at $w_r{=}10$ on the smaller cities would require capping $T_{\max}$ at $\sim 1$\,h or relaxing the load-factor floor.

\paragraph{Density scaling ($\Lambda$).} Three regimes emerge as $\Lambda$ rises from 50 to 300 at $w_r{=}1$. (i) Mohring (decreasing user time): SP-Lit drops 20--54\,\%, VCC drops 10--17\,\%, TP-Lit roughly flat. (ii) Anti-Mohring (increasing user time): Joint-Nom/DRO rise 19--34\,\% across all four cities because their Newton $T^*$ \emph{lengthens} with $\Lambda$ to keep tour cost competitive (a CA-literature prediction); OD on Pittsburgh ($+39\%$) and Manhattan ($+144\%$) saturates fleet queueing. (iii) $\Lambda$-insensitive: Multi-FR ($\pm 1\%$ everywhere) because additional demand fills the same already-scheduled trips at the locked $T_{\max}$ headway.

\paragraph{Where each baseline wins.} The cross-city cloud localises the operational choice. Multi-FR is the cheap-provider option (provider $20$--$30$, user time $\sim 1.6$\,h), tolerable only when riders accept long headways (rural or low-demand microtransit). Joint-Nom/DRO and SP-Lit are the provider-balanced option (provider 40--135, user time 0.4--0.9\,h at $w_r{=}1$); the DRO premium over Nom is small ($\sim 0$ to $15\%$) on population-weighted demand. TP-Lit, VCC, and OD are the user-time-optimised option (provider 90--1{,}220, user time 0.07--0.38\,h); cost rises super-linearly with topology irregularity, so the deployment city sharply affects the recommended choice.

\begin{figure}[t]
\centering
\includegraphics[width=\linewidth]{figures/external baselines/designs_four_cities}
\caption{Service designs across the four cities at $\Lambda{=}150$, $w_r{=}1$, with $K{=}3$ fixed for visualisation. Rows are cities (Pittsburgh, Chicago, Manhattan, LA); columns are baselines (Joint-Nom, Joint-DRO, SP-Lit, Multi-FR, TP-Lit, VCC, OD). Route polylines follow road shortest paths. Multi-FR shows selected reference loops; TP-Lit shows three dispatch waves; VCC shows walk halos; OD shows depot plus sample radial trips. Only Joint-Nom and Joint-DRO optimise the depot.}
\label{fig:four-city-designs}
\end{figure}

\paragraph{Implication.} The synthetic-grid conclusions (\S\ref{sec:cloud-analysis}) survive on real cities: the qualitative ranking is topology-robust, and the architectural distinctions remain visible in both the cloud and the design maps. The absolute position of the user-time regime varies with topology by a factor of 2 to 3 across cities, so deployment recommendations should be made on the actual target city, not extrapolated from a synthetic instance. Multi-FR remains the main outlier: its paper-faithful coverage reward yields a competitive provider cost but a long-headway design that pins user time near $T_{\max}/2$ regardless of regime. Reducing $T_{\max}$ or relaxing the load-factor floor would let Multi-FR enter a more user-favorable region of the Pareto plane, at the cost of departing from the paper's calibration.
 where needed.
% Requires: amsmath, booktabs, array.
% Figure paths assume Overleaf layout: figures/external baselines/<name>
% =============================================================================

\section{External Baselines}\label{sec:external-baselines}

We benchmark our service-design framework against seven external baselines drawn from three families: (i) partition with continuous approximation (\textbf{Joint-Nom}, \textbf{Joint-DRO}, \textbf{SP-Lit (Carlsson)}); (ii) routing-based (\textbf{Multi-FR (Jacquillat)}, \textbf{TP-Lit (Liu)}); (iii) operational extremes (\textbf{VCC}, \textbf{OD/fully-on-demand}). Each baseline sweeps its own structural hyperparameters internally per regime point and picks the combination minimising the aggregate cost $C_{\mathrm{prov}} + w_r \cdot \bar{u} \cdot \Lambda$. The comparison is stress-tested at three tiers: \S\ref{sec:cloud-analysis} synthetic $5{\times}5$ grid; \S\ref{sec:real-cloud} a 25-block-group Pittsburgh slice with real road distances and ACS population demand; \S\ref{sec:multi-city} four U.S.\ cities spanning topology archetypes (Pittsburgh, Chicago, Manhattan, LA). To make the regime response readable, each cloud is presented in two faceted views: one figure with $\Lambda$ fixed per panel and $w_r$ varying within the panel (shows how each baseline responds to the user-time weight at a given demand level), and a companion figure with $w_r$ fixed per panel and $\Lambda$ varying within the panel (shows how each baseline responds to demand at a given user-time weight).

\subsection{Analysis on the synthetic instance}\label{sec:cloud-analysis}

\begin{figure}[t]
\centering
\includegraphics[width=0.92\linewidth]{figures/external baselines/cloud_synthetic}
\caption{Performance cloud on the $5{\times}5$ synthetic instance ($\Lambda \in \{50, 150, 300\}$\,pax/h, $w_r \in \{1, 3, 10\}$, $w_v{=}1$, $c_f{=}4.0$\,\$/veh-h). Marker size $\propto \Lambda$, fill shade $\propto w_r$ (lighter $=$ larger $w_r$), shaded hull is each baseline's spread over the $3{\times}3$ regime grid.}
\label{fig:cloud}
\end{figure}

\begin{figure}[t]
\centering
\includegraphics[width=\linewidth]{figures/external baselines/cloud_synthetic_by_lambda}
\caption{Synthetic cloud, fixed $\Lambda$ per panel, $w_r$ varying within panel (fill shade encodes $w_r$, lighter $=$ smaller $w_r$). Lines connect the $w_r$ sweep for each baseline. Baselines whose design is genuinely $w_r$-insensitive at a given $\Lambda$ (TP-Lit at every $\Lambda$, VCC and OD at low $\Lambda$) have all three $w_r$ rows at identical (provider, user) coordinates; markers are then jittered slightly in $\log x$ so all three remain visible.}
\label{fig:cloud-by-lambda}
\end{figure}

\begin{figure}[t]
\centering
\includegraphics[width=\linewidth]{figures/external baselines/cloud_synthetic_by_wr}
\caption{Synthetic cloud, fixed $w_r$ per panel, $\Lambda$ varying within panel (marker size $\propto \Lambda$). Lines connect the $\Lambda$ sweep for each baseline.}
\label{fig:cloud-by-wr}
\end{figure}

Figure~\ref{fig:cloud} plots each baseline's provider cost rate (log scale) against average user time per rider across the full $3 \times 3$ regime grid. Figures~\ref{fig:cloud-by-lambda} and~\ref{fig:cloud-by-wr} decompose the same data by holding one regime parameter fixed per panel. The by-$\Lambda$ view exposes each method's $w_r$ trajectory: partition methods (Joint-Nom, Joint-DRO, SP-Lit) sweep from cheap-provider/long-wait at $w_r{=}1$ to expensive-provider/short-wait at $w_r{=}10$, while TP-Lit, VCC, and OD slide only along the provider axis at roughly fixed user time, and Multi-FR barely moves. The by-$w_r$ view exposes each method's $\Lambda$ trajectory: Multi-FR shows the strongest Mohring effect (user time drops from $\approx 1.07$\,h at $\Lambda{=}50$ to $\approx 0.56$\,h at $\Lambda{=}300$, $w_r{=}1$), partition methods are weakly anti-Mohring (Joint's $T^*$ lengthens with $\Lambda$), and OD/VCC have nearly flat user time (the fleet scales up with $\Lambda$ to absorb demand, so all $\Lambda$ variation appears horizontally on the provider axis).

\subsubsection{Per-baseline sweep grids}

\begin{center}
\small
\renewcommand{\arraystretch}{1.15}
\begin{tabular}{lll}
\toprule
Baseline & Swept hyperparameters & Selection \\
\midrule
Joint-Nom  & $K \in \{2,\dots,6\}$ (random-search partition) & best $K$ by aggregate cost \\
Joint-DRO  & $K \in \{2,\dots,6\}$, $\varepsilon = 0.5$ fixed & best $K$ by aggregate cost \\
SP-Lit     & $K \in \{2,\dots,6\}$ $\times$ $T_{\mathrm{fixed}} \in \{0.1, 0.2, 0.3, 0.5, 0.8, 1.0\}$ & best $(K, T)$ \\
Multi-FR   & (fleet, $\delta$) $\in$ $\{(3,0), (5,0), (8,0), (12,0)\}$, internal $T$-grid & best combo \\
TP-Lit     & $\mathtt{max\_fleet} \in \{20, 50, 100\}$; 25 candidate stops & best fleet \\
VCC        & fleet $\in \{5,10,20,30,50\} \times \delta_w \in \{0.1, 0.3, 0.5\}$\,km & best (fleet, $\delta_w$) \\
OD         & fleet $\in \{5,10,15,20,30,50,80\}$ with queueing & best fleet \\
\bottomrule
\end{tabular}
\end{center}

Joint-DRO holds $\varepsilon{=}0.5$ fixed because it is a robustness knob, not a provider/user Pareto knob.

\subsubsection{Joint-Nom and Joint-DRO}

Joint methods occupy the left band of the cloud (provider cost 28 to 136), spanning a wide user-time range (0.11 to 0.78\,h). Two design knobs respond to the regime: the number of districts $K$ (swept externally over $\{2,\dots,6\}$ by the harness, best $K$ picked per regime point by aggregate cost) and the per-district dispatch interval $T^*$ (computed internally by Newton's method for each candidate $K$).

\begin{center}
\small
\begin{tabular}{cccccc}
\toprule
$\Lambda$ & $w_r$ & $K$ & $T^*$ (h) & User time (h) & Provider cost \\
\midrule
 50 &  1 & 4 & 0.88 & 0.60 &  28 \\
 50 & 10 & 5 & 0.13 & 0.12 &  50 \\
150 &  1 & 6 & 0.99 & 0.70 &  38 \\
150 & 10 & 6 & 0.12 & 0.16 &  95 \\
300 &  1 & 6 & 1.00 & 0.78 &  43 \\
300 & 10 & 6 & 0.13 & 0.19 & 124 \\
\bottomrule
\end{tabular}
\end{center}

$T^*$ is highly $w_r$-responsive: at $w_r{=}1$ the optimiser values rider time cheaply and picks $T^* \approx 0.9$ to $1.0$\,h (heavy batching), while at $w_r{=}10$ it drops to $\approx 0.13$\,h (near real-time dispatch) at 3 to 4$\times$ the provider cost. $K$ increases with both $w_r$ and $\Lambda$ (from $K{=}4$ at the cheap end to $K{=}6$ at the expensive end). User time barely responds to $\Lambda$ at low $w_r$ because the Newton $T^*$ formula uses a per-rider wait term $w_r \cdot T$ rather than aggregate $w_r \cdot T \cdot \Lambda$, and the $K$ sweep partially compensates by adding districts at higher $\Lambda$.

Joint-DRO separates from Joint-Nom only at high $\Lambda$ and $w_r$: at $(\Lambda{=}300, w_r{=}10)$, DRO picks fleet 17 / provider cost 136 versus Nom's fleet 15 / provider cost 124 ($\sim 10\%$ premium). Under uniform demand the DRO premium stays modest; non-uniform demand amplifies it.

\subsubsection{Multi-FR (Jacquillat)}

Multi-FR occupies the mid-left region and is the only method with $\Lambda$-decreasing user time on the synthetic grid (the Mohring effect): more riders justify shorter dispatch headways on the same fixed-route backbone, so $T^*$ scales roughly as $1/\sqrt{\Lambda}$. Fleet 3 wins at every synthetic regime point; fleet 5 is available but rarely preferred.

\begin{center}
\small
\begin{tabular}{cccccc}
\toprule
$\Lambda$ & $w_r$ & $T$ (h) & Wait (h) & User time (h) & Provider cost \\
\midrule
 50 &  1 & 1.75 & 0.83 & 1.04 & 30 \\
150 &  1 & 0.81 & 0.40 & 0.62 & 40 \\
300 &  1 & 0.69 & 0.35 & 0.56 & 44 \\
300 & 10 & 0.38 & 0.19 & 0.39 & 69 \\
\bottomrule
\end{tabular}
\end{center}

Multi-FR is one of the few baselines with non-zero walk time ($\approx 0.07$ to $0.12$\,h on the grid), reflecting the rider-to-checkpoint distance. The evaluator assigns each sampled rider to the nearest realised service point globally rather than to a partition-assigned route.

\subsubsection{SP-Lit (Carlsson)}

SP-Lit sits in the centre-right of the cloud, shifted right of Joint despite the same partition-plus-BHH structure. At every regime point its provider cost is 1.6 to 3.1$\times$ Joint-Nom's. The gap is a partition-quality difference: SP-Lit's deterministic ham-sandwich bisection balances the continuous measures $\mu_1$ (spatial) and $\mu_2$ (radial workload) but does not directly minimise per-district cost; Joint's random search with local improvement finds lower-$\alpha_i$ partitions for the same $K$.

\subsubsection{TP-Lit (Liu)}\label{sec:tplit-cloud}

TP-Lit occupies a narrow horizontal band at user time $\approx 0.20$ to $0.23$\,h with provider cost 233 to 244. On the synthetic instance the design is essentially $w_r$-insensitive: the ADP+VRP method ran at its library defaults ($T{=}0.25$\,h, fleet automatically sized by APT to 22), the three $w_r$ rows at each $\Lambda$ are bit-identical, and the three visible markers per $\Lambda$ panel sit on the same coordinate (small log-$x$ jitter applied for visibility). The real-city sweeps (\S\ref{sec:real-cloud}, \S\ref{sec:multi-city}) do add an external $T_{\mathrm{fixed}} \in \{0.15, 0.5\}$\,h $\times$ $\mathtt{max\_fleet} \in \{30, 100\}$ wrapper sweep and pick the best 4-combo by aggregate cost, but the chosen combo is the same at every $w_r$ at most regime points (e.g.\ all three $w_r$ at $\Lambda{=}150$ Pittsburgh pick $T{=}0.5$, $\mathtt{max\_fleet}{=}30$); a $T{=}0.15$ switch only occurs at $w_r{=}10$ on $\Lambda \in \{50, 300\}$. The reason is structural: TP-Lit's inner objective (ADP delivery-time minimisation) does not contain $w_r$, so only the outer wrapper's aggregate-cost selection responds to $w_r$ at all, and that response is quantised by the 4-combo grid. Maintaining a sizeable fleet for 25 candidate stops makes the fleet-cost term alone ($c_f \times \mathtt{fleet}$) account for $\sim 35\%$ of provider cost on the synthetic.

\subsubsection{VCC and OD}

VCC and OD form the operational-extremes panel (Fig.~\ref{fig:cloud-split}b). Both achieve user time below 0.16\,h on the synthetic grid but at provider cost 1 to 2 orders of magnitude above Joint-Nom. VCC sweeps fleet ($\{5,10,20,30,50\}$) and walk tolerance $\delta_w \in \{0.1,0.3,0.5\}$\,km; the optimiser picks larger fleets at higher $\Lambda$ to keep user time constant ($\approx 0.15$\,h), producing a horizontal trajectory in the cloud. OD serves each request individually (depot $\to$ destination $\to$ depot) and is the cost ceiling on provider expense and the cost floor on user time: every other baseline is interpreted relative to OD's extreme. At low $w_r$, OD tolerates moderate queueing (fleet 5--15); at $w_r{=}10$ it upgrades to 50 vehicles at $\Lambda{=}300$.

\subsubsection{Cross-method synthesis}

No single method dominates across the regime grid. The approximate Pareto frontier runs from Joint-Nom/DRO (cheapest provider, highest user time at low $w_r$) through Multi-FR (Mohring-aided at high $\Lambda$) to VCC and OD (lowest user time, highest provider cost). SP-Lit and TP-Lit are dominated on the frontier: SP-Lit by Joint-Nom (same structure, better partitions) and TP-Lit by VCC (lower user time at comparable provider cost). The clearest architectural distinction is the response to $\Lambda$: Multi-FR's user time \emph{decreases} (Mohring); Joint/SP-Lit barely respond (the $K$ sweep absorbs the demand change); VCC/OD are nearly constant in user time (fleet scales with $\Lambda$); TP-Lit is perfectly flat (its design is $\Lambda$-insensitive in this range).

\subsection{Pittsburgh real-data validation}\label{sec:real-cloud}

\begin{figure}[t]
\centering
\includegraphics[width=\linewidth]{figures/external baselines/cloud_pittsburgh}
\caption{Pittsburgh real-data cloud, 25 census block groups in Allegheny County, road-network shortest paths everywhere (optimisation and evaluation), ACS 2023 population-weighted demand. Same regime grid and same axes as Fig.~\ref{fig:cloud}.}
\label{fig:cloud-real}
\end{figure}

\begin{figure}[t]
\centering
\includegraphics[width=\linewidth]{figures/external baselines/cloud_pittsburgh_by_lambda}
\caption{Pittsburgh cloud, fixed $\Lambda$ per panel, $w_r$ varying.}
\label{fig:cloud-real-by-lambda}
\end{figure}

\begin{figure}[t]
\centering
\includegraphics[width=\linewidth]{figures/external baselines/cloud_pittsburgh_by_wr}
\caption{Pittsburgh cloud, fixed $w_r$ per panel, $\Lambda$ varying.}
\label{fig:cloud-real-by-wr}
\end{figure}

We re-run the same regime grid on 25 census block groups around the centroid of HCT's Pittsburgh service area. Demand is ACS-population-weighted ($p_b \in [0.014, 0.062]$ versus the synthetic $p_b{=}1/25$); distances are OSMnx road-network shortest paths (2{,}931 nodes, 7{,}519 edges) everywhere a baseline supports them (Joint-Nom, Joint-DRO, SP-Lit in optimisation; all baselines in evaluation).

\paragraph{Headline numbers at $\Lambda{=}150$, $w_r{=}1$.}
\begin{center}
\small
\begin{tabular}{lcc}
\toprule
Baseline & Provider cost & User time (h) \\
\midrule
Multi-FR (Jacquillat)  &  29.5 & 1.78 \\
Joint-Nom              &  58.6 & 0.92 \\
Joint-DRO              &  59.5 & 0.91 \\
SP-Lit (Carlsson)      & 134.9 & 0.70 \\
TP-Lit (Liu)           & 216.2 & 0.38 \\
VCC                    & 718.1 & 0.21 \\
OD                     & 1220.2 & 0.23 \\
\bottomrule
\end{tabular}
\end{center}

\paragraph{Findings.} Three patterns from the synthetic analysis carry over and one new one emerges. (i) Provider costs scale up across the board because real road distances are longer than Euclidean (typical road-to-Euclidean factor 1.3 to 1.5 in Pittsburgh's grid); OD's provider cost more than triples, while partition methods absorb most of the increase via their road-aware in-district TSPs. (ii) The Pareto ranking is preserved: same ordering Joint $\to$ SP-Lit $\to$ TP-Lit $\to$ VCC $\to$ OD on provider cost, with the reverse on user time. (iii) Joint's $K$ choice tracks regime monotonically ($K{=}4$ at $\Lambda{=}50, w_r{=}1$, $K{=}6$ at $\Lambda{=}300, w_r{=}10$). (iv) Non-uniform demand sharpens the DRO-vs-Nom premium: at $(\Lambda{=}300, w_r{=}10)$, DRO costs $\sim 25\%$ more than Nom (versus $\sim 10\%$ on uniform synthetic), matching the prediction that DRO's worst-case partition diverges more aggressively when demand is spatially non-uniform.

\begin{figure}[t]
\centering
\includegraphics[width=\linewidth]{figures/external baselines/designs_pittsburgh}
\caption{Pittsburgh service designs at $\Lambda{=}150$, $w_r{=}1$, with $K{=}3$ fixed for visualisation (the cloud sweep picks $K$ per regime). Route polylines follow road shortest paths. Partition methods (cols 1 to 3) show contiguous district fills with exact-TSP depot tours. Multi-FR (col 4) shows selected reference loops. TP-Lit (col 5) shows three independent dispatch waves (blues, reds, greens) with realised VRP routes per vehicle. VCC (col 6) shows walk-radius halos ($\delta_w{=}0.3$\,km). OD (col 7) shows depot plus sample on-demand radial trips.}
\label{fig:designs-map}
\end{figure}

The map overlay confirms that the architectural distinctions are also visible in geometry: Joint-Nom and Joint-DRO produce contiguous random-search partitions with a Joint-Nom-optimised depot; SP-Lit's ham-sandwich partition is more axis-aligned; Multi-FR's selected reference loops radiate from the centroid depot; VCC overlays per-block walk halos; OD shows pure depot-anchored trips.

\subsection{Cross-topology validation across four U.S.\ cities}\label{sec:multi-city}

\begin{figure}[t]
\centering
\includegraphics[width=\linewidth]{figures/external baselines/cloud_four_cities}
\caption{Cross-city baseline performance clouds on the shared regime grid. Each panel is one city; markers and shading as in Fig.~\ref{fig:cloud}.}
\label{fig:four-city-cloud}
\end{figure}

\begin{figure}[t]
\centering
\includegraphics[width=\linewidth]{figures/external baselines/cloud_four_cities_by_lambda}
\caption{Cross-city clouds, fixed $\Lambda$ per column, $w_r$ varying within panel. Rows are cities.}
\label{fig:four-city-cloud-by-lambda}
\end{figure}

\begin{figure}[t]
\centering
\includegraphics[width=\linewidth]{figures/external baselines/cloud_four_cities_by_wr}
\caption{Cross-city clouds, fixed $w_r$ per column, $\Lambda$ varying within panel. Rows are cities.}
\label{fig:four-city-cloud-by-wr}
\end{figure}

We re-run the regime sweep on four cities chosen for road-topology contrast: Pittsburgh (irregular hilly, river-cut), Chicago Loop (strict cardinal grid), Manhattan Midtown (avenue-and-street grid), and Los Angeles DTLA (freeway-fragmented grid). Each city uses real TIGER 2023 block groups, ACS-population demand, and full road-network shortest paths in both optimisation and evaluation.

\paragraph{Service-region geometry.} BG count is held at 25 (50 for Manhattan, see below); TIGER block groups are population-sized, so denser cities pack the same count into smaller polygons.

\begin{center}
\small
\begin{tabular}{lcrrrr}
\toprule
City & BGs & Area (km$^2$) & BBox W$\times$H (km) & Diag (km) & Density (pop/km$^2$) \\
\midrule
Pittsburgh & 25 & 24.5 & $7.4 \times 7.2$ & 10.3 &  1{,}003 \\
Chicago    & 25 &  4.6 & $2.3 \times 3.5$ &  4.1 & 13{,}937 \\
Manhattan  & 50 &  2.4 & $2.0 \times 2.5$ &  3.2 & 18{,}891 \\
LA         & 25 &  4.3 & $2.3 \times 2.9$ &  3.7 &  8{,}755 \\
\bottomrule
\end{tabular}
\end{center}

Pittsburgh's 25 BGs cover roughly $10{\times}$ more area than Manhattan's 50, so cross-city cost comparisons reflect both topology and area. Manhattan is enlarged from the default 25 BGs ($1.1$\,km$^2$, too small for FR-Detour's candidate-line library to produce feasible trips) to 50 BGs ($2.4$\,km$^2$); the other three cities are kept at 25 BGs.

\paragraph{Headline numbers at $\Lambda{=}150$, $w_r{=}1$.}
\begin{center}
\small
\begin{tabular}{lrrrr}
\toprule
Baseline & Pittsburgh & Chicago & Manhattan & LA \\
\midrule
\multicolumn{5}{l}{\emph{Provider cost}} \\
Joint-Nom              &   58.6 &   47.6 &   47.9 &   35.5 \\
Joint-DRO              &   59.5 &   45.8 &   54.6 &   36.7 \\
SP-Lit (Carlsson)      &  134.9 &   81.1 &   73.7 &   58.5 \\
Multi-FR (Jacquillat)  &   29.5 &   21.4 &   20.5 &   20.5 \\
TP-Lit (Liu)           &  216.2 &  103.1 &  104.7 &   91.5 \\
VCC                    &  718.1 &  241.9 &  199.7 &  193.3 \\
OD                     & 1220.2 &  476.4 &  417.5 &  343.8 \\
\midrule
\multicolumn{5}{l}{\emph{User time (h per rider)}} \\
Joint-Nom              & 0.92 & 0.64 & 0.58 & 0.56 \\
Joint-DRO              & 0.91 & 0.68 & 0.60 & 0.56 \\
SP-Lit (Carlsson)      & 0.70 & 0.41 & 0.39 & 0.41 \\
Multi-FR (Jacquillat)  & 1.78 & 1.61 & 1.58 & 1.59 \\
TP-Lit (Liu)           & 0.38 & 0.31 & 0.31 & 0.30 \\
VCC                    & 0.21 & 0.12 & 0.11 & 0.09 \\
OD                     & 0.23 & 0.07 & 0.18 & 0.08 \\
\bottomrule
\end{tabular}
\end{center}

\paragraph{Multi-FR calibration.} Multi-FR uses the paper-faithful coverage reward $M = 10{,}000 \cdot w_r$ (Jacquillat et al.\ Table EC.4). With this calibration the master picks fleet $\in \{4, 6, 9, 12\}$ across the regime grid (versus the degenerate fleet $=1$ under our earlier $M = 10 \cdot w_r$), and walk time drops from 0.4 to 0.5\,h down to $\sim 0.07$\,h. However, the master still selects the longest scheduled headway in our $T$-grid ($T_{\max}{=}3$\,h, realised wait $\approx 1.5$\,h) at every regime point because $M{=}10^4 \cdot w_r$ dwarfs the per-passenger wait term $w_r \cdot T/2 \cdot \mathrm{demand}$ by $\sim 6{,}700{\times}$. Multi-FR therefore lands at a many-trips, full-coverage, long-headway design point on real cities. At $w_r{=}10$ on Manhattan and LA the master sometimes flips to an empty-solution corner of the feasible set; see the discussion below.

\paragraph{Cross-method ranking and Pareto positions.} The provider-cost ordering Multi-FR $<$ Joint $<$ SP-Lit $<$ TP-Lit $<$ VCC $<$ OD holds at every $(\Lambda, w_r)$ point on all four cities; the user-time ordering is essentially the reverse. Multi-FR consistently occupies the cheapest-provider corner (provider 20--30, user time 1.6--1.8\,h); OD the opposite (provider 344--1{,}220, user time 0.07--0.23\,h). The Pareto frontier is dense: at $(\Lambda{=}150, w_r{=}1)$, five to seven of the seven baselines are non-dominated on each city. The architectural distinctions between baselines correspond to genuinely different operating points on the same Pareto plane.

\paragraph{Topology controls absolute cost.} Pittsburgh (25 BGs, 24.5\,km$^2$, irregular) inflates absolute provider cost relative to Manhattan (50 BGs, 2.4\,km$^2$, dense grid): Joint-Nom $59$ vs $48$ ($1.2{\times}$), SP-Lit $135$ vs $74$ ($1.8{\times}$), TP-Lit $216$ vs $105$ ($2.1{\times}$), OD $1{,}220$ vs $418$ ($2.9{\times}$). Methods that pay distance per request (OD, VCC) amplify the geometric handicap; partition methods absorb most of it via a shared in-district TSP. Multi-FR is essentially topology-invariant in provider cost (clustering at 20--30 across all four cities) because its long-headway design is dispatch-frequency-bound, not distance-bound. The OD-to-Joint multiplier rises with topology complexity: $9{\times}$ on Manhattan, 10--13$\times$ on Chicago and LA, $21{\times}$ on Pittsburgh.

\paragraph{Regime responsiveness ($w_r$).} As $w_r$ rises from 1 to 10 at $\Lambda{=}150$, partition methods (Joint-Nom/DRO, SP-Lit) drop user time by 66 to 77\,\% on every city. TP-Lit and OD also respond (23 to 69\,\%); VCC is weakly responsive (10 to 13\,\%). Multi-FR is essentially $w_r$-insensitive ($-1\%$ to $-4\%$ on Pittsburgh and Chicago), and degenerates to fleet$=0$ on Manhattan and LA at $w_r{=}10$. This is a known artifact of the calibrated $M = 10{,}000 \cdot w_r$ objective: at high $w_r$ both the per-passenger reward and the wait term scale up, and the master can satisfy load-factor and trip-budget constraints either by selecting $\sim 4$ trips or by selecting zero; on the smaller cities the empty corner sometimes wins. A non-degenerate Multi-FR at $w_r{=}10$ on the smaller cities would require capping $T_{\max}$ at $\sim 1$\,h or relaxing the load-factor floor.

\paragraph{Density scaling ($\Lambda$).} Three regimes emerge as $\Lambda$ rises from 50 to 300 at $w_r{=}1$. (i) Mohring (decreasing user time): SP-Lit drops 20--54\,\%, VCC drops 10--17\,\%, TP-Lit roughly flat. (ii) Anti-Mohring (increasing user time): Joint-Nom/DRO rise 19--34\,\% across all four cities because their Newton $T^*$ \emph{lengthens} with $\Lambda$ to keep tour cost competitive (a CA-literature prediction); OD on Pittsburgh ($+39\%$) and Manhattan ($+144\%$) saturates fleet queueing. (iii) $\Lambda$-insensitive: Multi-FR ($\pm 1\%$ everywhere) because additional demand fills the same already-scheduled trips at the locked $T_{\max}$ headway.

\paragraph{Where each baseline wins.} The cross-city cloud localises the operational choice. Multi-FR is the cheap-provider option (provider $20$--$30$, user time $\sim 1.6$\,h), tolerable only when riders accept long headways (rural or low-demand microtransit). Joint-Nom/DRO and SP-Lit are the provider-balanced option (provider 40--135, user time 0.4--0.9\,h at $w_r{=}1$); the DRO premium over Nom is small ($\sim 0$ to $15\%$) on population-weighted demand. TP-Lit, VCC, and OD are the user-time-optimised option (provider 90--1{,}220, user time 0.07--0.38\,h); cost rises super-linearly with topology irregularity, so the deployment city sharply affects the recommended choice.

\begin{figure}[t]
\centering
\includegraphics[width=\linewidth]{figures/external baselines/designs_four_cities}
\caption{Service designs across the four cities at $\Lambda{=}150$, $w_r{=}1$, with $K{=}3$ fixed for visualisation. Rows are cities (Pittsburgh, Chicago, Manhattan, LA); columns are baselines (Joint-Nom, Joint-DRO, SP-Lit, Multi-FR, TP-Lit, VCC, OD). Route polylines follow road shortest paths. Multi-FR shows selected reference loops; TP-Lit shows three dispatch waves; VCC shows walk halos; OD shows depot plus sample radial trips. Only Joint-Nom and Joint-DRO optimise the depot.}
\label{fig:four-city-designs}
\end{figure}

\paragraph{Implication.} The synthetic-grid conclusions (\S\ref{sec:cloud-analysis}) survive on real cities: the qualitative ranking is topology-robust, and the architectural distinctions remain visible in both the cloud and the design maps. The absolute position of the user-time regime varies with topology by a factor of 2 to 3 across cities, so deployment recommendations should be made on the actual target city, not extrapolated from a synthetic instance. Multi-FR remains the main outlier: its paper-faithful coverage reward yields a competitive provider cost but a long-headway design that pins user time near $T_{\max}/2$ regardless of regime. Reducing $T_{\max}$ or relaxing the load-factor floor would let Multi-FR enter a more user-favorable region of the Pareto plane, at the cost of departing from the paper's calibration.
